% Options for packages loaded elsewhere
\PassOptionsToPackage{unicode}{hyperref}
\PassOptionsToPackage{hyphens}{url}
%
\documentclass[
  a4paper,
  10pt]{article}
\usepackage{amsmath,amssymb}
\usepackage{iftex}
\ifPDFTeX
  \usepackage[T1]{fontenc}
  \usepackage[utf8]{inputenc}
  \usepackage{textcomp} % provide euro and other symbols
\else % if luatex or xetex
  \usepackage{unicode-math} % this also loads fontspec
  \defaultfontfeatures{Scale=MatchLowercase}
  \defaultfontfeatures[\rmfamily]{Ligatures=TeX,Scale=1}
\fi
\usepackage{lmodern}
\ifPDFTeX\else
  % xetex/luatex font selection
  \setmainfont[]{Lato}
  \setsansfont[]{Lato}
  \setmonofont[]{DejaVu Sans Mono}
  \setmathfont[]{STIXMath-Regular.otf}
\fi
% Use upquote if available, for straight quotes in verbatim environments
\IfFileExists{upquote.sty}{\usepackage{upquote}}{}
\IfFileExists{microtype.sty}{% use microtype if available
  \usepackage[]{microtype}
  \UseMicrotypeSet[protrusion]{basicmath} % disable protrusion for tt fonts
}{}
\makeatletter
\@ifundefined{KOMAClassName}{% if non-KOMA class
  \IfFileExists{parskip.sty}{%
    \usepackage{parskip}
  }{% else
    \setlength{\parindent}{0pt}
    \setlength{\parskip}{6pt plus 2pt minus 1pt}}
}{% if KOMA class
  \KOMAoptions{parskip=half}}
\makeatother
\usepackage{xcolor}
\usepackage{longtable,booktabs,array}
\usepackage{calc} % for calculating minipage widths
% Correct order of tables after \paragraph or \subparagraph
\usepackage{etoolbox}
\makeatletter
\patchcmd\longtable{\par}{\if@noskipsec\mbox{}\fi\par}{}{}
\makeatother
% Allow footnotes in longtable head/foot
\IfFileExists{footnotehyper.sty}{\usepackage{footnotehyper}}{\usepackage{footnote}}
\makesavenoteenv{longtable}
\usepackage{graphicx}
\makeatletter
\def\maxwidth{\ifdim\Gin@nat@width>\linewidth\linewidth\else\Gin@nat@width\fi}
\def\maxheight{\ifdim\Gin@nat@height>\textheight\textheight\else\Gin@nat@height\fi}
\makeatother
% Scale images if necessary, so that they will not overflow the page
% margins by default, and it is still possible to overwrite the defaults
% using explicit options in \includegraphics[width, height, ...]{}
\setkeys{Gin}{width=\maxwidth,height=\maxheight,keepaspectratio}
% Set default figure placement to htbp
\makeatletter
\def\fps@figure{htbp}
\makeatother
\setlength{\emergencystretch}{3em} % prevent overfull lines
\providecommand{\tightlist}{%
  \setlength{\itemsep}{0pt}\setlength{\parskip}{0pt}}
\setcounter{secnumdepth}{-\maxdimen} % remove section numbering
\ifLuaTeX
\usepackage[bidi=basic]{babel}
\else
\usepackage[bidi=default]{babel}
\fi
\babelprovide[main,import]{ngerman}
\ifPDFTeX
\else
\babelfont{rm}[]{Lato}
\fi
% get rid of language-specific shorthands (see #6817):
\let\LanguageShortHands\languageshorthands
\def\languageshorthands#1{}

% --- Dokumentdesign -------------------------------------------------------
\usepackage{microtype}
\usepackage{xcolor}
\definecolor{Brand}{HTML}{167B88}
\definecolor{Ink}{HTML}{14304F}
\definecolor{InkSoft}{HTML}{2B4F77}
\definecolor{Muted}{HTML}{5C636A}
\definecolor{Rule}{HTML}{DDE3E8}
\definecolor{Wash}{HTML}{F4F7F8}
\color{Ink}

\usepackage[a4paper,left=19mm,right=19mm,top=24mm,bottom=21mm,
            headheight=18pt,headsep=8mm,footskip=12mm]{geometry}
\usepackage{setspace}
\setstretch{1.08}
\setlength{\parindent}{0pt}
\setlength{\parskip}{5.5pt plus 1pt minus 1pt}
\emergencystretch=2.5em

\usepackage{titlesec}
\titleformat{\section}
  {\sffamily\bfseries\fontsize{20}{24}\selectfont\color{Ink}}
  {}{0pt}{}
\titleformat{\subsection}
  {\sffamily\bfseries\fontsize{13.2}{16}\selectfont\color{Brand}}
  {}{0pt}{}
\titleformat{\subsubsection}
  {\sffamily\bfseries\fontsize{10.7}{13.5}\selectfont\color{InkSoft}}
  {}{0pt}{}
\titlespacing*{\section}{0pt}{0pt}{13pt}
\titlespacing*{\subsection}{0pt}{17pt}{6pt}
\titlespacing*{\subsubsection}{0pt}{12pt}{4pt}
\newcommand{\sectionbreak}{\clearpage}

\usepackage{fancyhdr}
\pagestyle{fancy}
\fancyhf{}
\fancyhead[L]{\sffamily\fontsize{6.6}{8}\selectfont\color{Muted}Valyze · Dokumentation Cash-Flow-Model - Rechenmodell und Berechnungsvorschrift}
\fancyhead[R]{}
\fancyfoot[L]{\sffamily\fontsize{6.6}{8}\selectfont\color{Muted}\nouppercase{\leftmark}}
\fancyfoot[R]{\sffamily\fontsize{7.6}{9}\selectfont\color{Muted}Seite \thepage}
\renewcommand{\headrulewidth}{0.35pt}
\renewcommand{\footrulewidth}{0.35pt}
\renewcommand{\headrule}{\hbox to\headwidth{\color{Rule}\leaders\hrule height \headrulewidth\hfill}}
\renewcommand{\footrule}{\hbox to\headwidth{\color{Rule}\leaders\hrule height \footrulewidth\hfill}}
\renewcommand{\sectionmark}[1]{\markboth{#1}{}}
\fancypagestyle{plain}{%
  \fancyhf{}
  \fancyfoot[C]{\sffamily\fontsize{7.6}{9}\selectfont\color{Muted}Seite \thepage}
  \renewcommand{\headrulewidth}{0pt}
  \renewcommand{\footrulewidth}{0.35pt}
}

\usepackage{amsmath,amssymb}
\setlength{\abovedisplayskip}{10pt plus 2pt minus 2pt}
\setlength{\belowdisplayskip}{11pt plus 2pt minus 2pt}
\setlength{\abovedisplayshortskip}{8pt plus 2pt}
\setlength{\belowdisplayshortskip}{9pt plus 2pt}
\setlength{\jot}{6pt}
\allowdisplaybreaks[1]
\newcommand{\mathbox}[1]{%
  \begingroup\setlength{\fboxsep}{8pt}%
  \colorbox{Wash}{\ensuremath{\displaystyle #1}}\endgroup}

\usepackage{booktabs,longtable,array,tabularx,colortbl}
\arrayrulecolor{Rule}
\renewcommand{\arraystretch}{1.16}
\setlength{\tabcolsep}{5pt}
\usepackage{etoolbox}
\AtBeginEnvironment{longtable}{\small\rowcolors{2}{Wash}{white}}
\AtBeginEnvironment{table}{\small}

\usepackage{enumitem}
\setlist[itemize]{leftmargin=17pt,itemsep=2.5pt,topsep=4pt}
\setlist[enumerate]{leftmargin=19pt,itemsep=2.5pt,topsep=4pt}
\setlist[itemize,1]{label=\textcolor{Brand}{\small\textbullet}}

\usepackage[most]{tcolorbox}
\renewenvironment{quote}
  {\begin{tcolorbox}[
      enhanced,
      breakable,
      colback=Wash,
      colframe=Brand,
      boxrule=0pt,
      leftrule=2.2pt,
      arc=0pt,
      left=8pt,right=8pt,top=6pt,bottom=6pt,
      before skip=8pt,after skip=9pt
    ]\small\color{InkSoft}}
  {\end{tcolorbox}}

\definecolor{shadecolor}{HTML}{F4F7F8}
\usepackage{fvextra}
\fvset{fontsize=\scriptsize,breaklines=true,breakanywhere=true,
       frame=single,rulecolor=\color{Rule},framesep=5pt}

\usepackage{caption}
\captionsetup{font=small,labelfont={bf,color=InkSoft},textfont={color=Muted},skip=6pt}
\usepackage{graphicx}
\setkeys{Gin}{width=\linewidth,height=0.72\textheight,keepaspectratio}

\usepackage{hyperref}
\usepackage{xurl}
\urlstyle{tt}
\Urlmuskip=0mu plus 1mu
\hypersetup{
  pdftitle={Dokumentation Cash-Flow-Model - Rechenmodell und Berechnungsvorschrift},
  pdfauthor={Valyze},
  pdfsubject={Mathematische Spezifikation der Projektbewertung},
  colorlinks=true,
  linkcolor=Brand,
  urlcolor=Brand,
  citecolor=Brand,
  bookmarksopen=true,
  bookmarksnumbered=false
}

% --- Individuelles Deckblatt ---------------------------------------------
\makeatletter
\renewcommand{\maketitle}{%
  \begin{titlepage}
    \thispagestyle{empty}
    \vspace*{15mm}
    \includegraphics[width=76mm]{valyze_marke.png}\par
    \vspace{10mm}
    {\sffamily\bfseries\fontsize{25}{30}\selectfont\color{Ink}
      Dokumentation Cash-Flow-Model -\par}
    \vspace{5mm}
    {\sffamily\bfseries\fontsize{17}{21}\selectfont\color{Brand}
      Rechenmodell und Berechnungsvorschrift\par}
    \vspace{11mm}
    {\color{Brand}\rule{56mm}{2.2pt}\par}
    \vspace{11mm}
    {\sffamily\fontsize{12}{17}\selectfont\color{InkSoft}
      Mathematische Spezifikation der Projektbewertung und Gebotsanalyse\par}
    \vfill
    \begin{tcolorbox}[
      colback=Wash,colframe=Rule,boxrule=0.5pt,arc=1.5pt,
      left=9pt,right=9pt,top=8pt,bottom=8pt,width=\textwidth
    ]
      {\sffamily\fontsize{9.2}{13}\selectfont\color{InkSoft}
      Wirtschaftlichkeitsrechnung für PV-Projekte nach dem
      österreichischen EAG-Marktprämienmodell\\[4pt]
      Stand: 16.08.2026 · erzeugt aus
      \nolinkurl{docs/rechenmodell/rechenmodell.md}}
    \end{tcolorbox}
  \end{titlepage}
  \setcounter{page}{1}
}
\makeatother
\ifLuaTeX
  \usepackage{selnolig}  % disable illegal ligatures
\fi
\IfFileExists{bookmark.sty}{\usepackage{bookmark}}{\usepackage{hyperref}}
\IfFileExists{xurl.sty}{\usepackage{xurl}}{} % add URL line breaks if available
\urlstyle{same}
\hypersetup{
  pdftitle={Dokumentation Cash-Flow-Model - Rechenmodell und Berechnungsvorschrift},
  pdfauthor={Valyze},
  pdflang={de-DE},
  hidelinks,
  pdfcreator={LaTeX via pandoc}}

\title{Dokumentation Cash-Flow-Model - Rechenmodell und
Berechnungsvorschrift}
\author{Valyze}
\date{16.08.2026}

\begin{document}
\maketitle

{
\setcounter{tocdepth}{3}
\tableofcontents
}
\hypertarget{gegenstand-geltungsbereich-und-dokumentstruktur}{%
\section{1 Gegenstand, Geltungsbereich und
Dokumentstruktur}\label{gegenstand-geltungsbereich-und-dokumentstruktur}}

\hypertarget{zielsetzung}{%
\subsection{1.1 Zielsetzung}\label{zielsetzung}}

Dieses Dokument enthält die \textbf{formale Spezifikation der im
Cash-Flow-Model implementierten Wirtschaftlichkeitsrechnung}. Sämtliche
im Modell erzeugten Größen werden durch ihre mathematische Definition,
ihre Eingangsparameter und ihre Stellung innerhalb der Berechnungsfolge
beschrieben. Ziel ist eine vollständige fachliche Nachvollziehbarkeit
der Modellergebnisse unabhängig von der konkreten
Softwareimplementierung.

Die Darstellung ermöglicht insbesondere

\begin{itemize}
\tightlist
\item
  die Rekonstruktion des aufgelösten Parametersatzes aus
  projektbezogenen und globalen Eingaben (Kapitel 4),
\item
  die unabhängige Reproduktion der Berechnungsschritte von der
  Periodisierung bis zu den Bewertungskennzahlen (Kapitel 5 bis 12),
\item
  die numerische Plausibilisierung anhand eines vollständig
  spezifizierten Beispielprojekts (Kapitel 13),
\item
  die methodische Einordnung der Sensitivitäts-, Risiko- und
  Break-even-Analysen (Kapitel 14),
\item
  die Dokumentation des empirischen Modells der EAG-Ausschreibungen
  einschließlich Verteilungsannahmen, Parameterschätzung,
  Wettbewerbsmodell und Prognoseverfahren (Kapitel 15) sowie
\item
  die Zuordnung der mathematischen Vorschriften zu Implementierung und
  Tests (Kapitel 18).
\end{itemize}

\hypertarget{sachliche-abgrenzung}{%
\subsection{1.2 Sachliche Abgrenzung}\label{sachliche-abgrenzung}}

Gegenstand des Dokuments ist ausschließlich die fachliche und
mathematische Rechenlogik. Die Softwarearchitektur wird in
\href{../ARCHITECTURE.md}{docs/ARCHITECTURE.md}, die Bedienung der
Anwendung in der \href{../../README.md}{README} beschrieben. Steuer- und
energierechtliche Regelungen werden nur insoweit dargestellt, wie sie
als explizite Berechnungsvorschriften implementiert sind. Die
Dokumentation stellt daher keine rechtliche oder steuerliche Würdigung
dar.

\hypertarget{aufbau-und-darstellungskonventionen}{%
\subsection{1.3 Aufbau und
Darstellungskonventionen}\label{aufbau-und-darstellungskonventionen}}

Die Kapitel zu den einzelnen Rechenmodulen folgen einer einheitlichen
Struktur:

\begin{enumerate}
\def\labelenumi{\arabic{enumi}.}
\tightlist
\item
  \textbf{Zielgröße} -- Definition der im jeweiligen Modul ermittelten
  Größe.
\item
  \textbf{Eingangsgrößen} -- erforderliche Parameter und Ergebnisse
  vorgelagerter Module.
\item
  \textbf{Berechnungsvorschrift} -- mathematische Definition in der
  Reihenfolge der Implementierung.
\item
  \textbf{Methodische Erläuterung} -- Begründung, Sonderfälle und
  Randbedingungen.
\item
  \textbf{Ausgangsgrößen} -- erzeugte Variablen und Einheiten.
\item
  \textbf{Implementierungsreferenz} -- zugehöriges Modul und zugehörige
  Funktion.
\end{enumerate}

\begin{quote}
Modellannahmen und bewusste Vereinfachungen werden an der jeweiligen
Berechnungsstelle gekennzeichnet und in Kapitel 17 konsolidiert.
\end{quote}

\hypertarget{mathematische-gesamtstruktur-des-bewertungsmodells}{%
\section{2 Mathematische Gesamtstruktur des
Bewertungsmodells}\label{mathematische-gesamtstruktur-des-bewertungsmodells}}

\hypertarget{berechnungsarchitektur-und-datenfluss}{%
\subsection{2.1 Berechnungsarchitektur und
Datenfluss}\label{berechnungsarchitektur-und-datenfluss}}

Die Projektbewertung ist als gerichteter Datenfluss organisiert. Jedes
Modul verwendet ausschließlich Eingangsparameter oder Ergebnisse
vorgelagerter Module. Algebraische Zirkelbezüge bestehen nicht. Eine
zeitliche Rekursion tritt lediglich bei zustandsabhängigen Größen auf,
insbesondere beim steuerlichen Verlustvortrag und beim Darlehensstand.
Diese Rekursionen werden periodenweise vorwärts ausgewertet und
erfordern keine iterative Gleichungslösung.

\begin{figure}
\centering
\includegraphics{rechenweg.png}
\caption{Berechnungsarchitektur der Projektbewertung}
\end{figure}

\hypertarget{bewertungsfunktion-und-hierarchisches-gleichungssystem}{%
\subsection{2.2 Bewertungsfunktion und hierarchisches
Gleichungssystem}\label{bewertungsfunktion-und-hierarchisches-gleichungssystem}}

Die Bewertung erfolgt aus Sicht der Eigenkapitalgeber auf Grundlage der
periodischen Equity-Cashflows. Die mathematische Gesamtstruktur besteht
aus drei Ebenen:

\begin{enumerate}
\def\labelenumi{\arabic{enumi}.}
\tightlist
\item
  dem Nettobarwert als expliziter Bewertungsfunktion,
\item
  dem internen Zinsfuß als implizit definierter Nullstelle dieser
  Funktion und
\item
  der periodenspezifischen Herleitung der zugrunde liegenden Cashflows.
\end{enumerate}

Diese Hierarchie ist fachlich zweckmäßiger als eine vollständig
ausmultiplizierte Einzelformel: Eine solche Darstellung wäre zwar
algebraisch möglich, würde jedoch die Abhängigkeiten zwischen
physischen, marktlichen, betrieblichen, finanziellen und steuerlichen
Größen verdecken.

\hypertarget{barwertfunktion-als-mathematischer-ausgangspunkt}{%
\subsubsection{2.2.1 Barwertfunktion als mathematischer
Ausgangspunkt}\label{barwertfunktion-als-mathematischer-ausgangspunkt}}

Für einen vorgegebenen Diskontsatz \(r\) ist der taggenaue Nettobarwert
durch

\[
\operatorname{XNPV}(r)
  = \sum_{t=0}^{N}
    \frac{\mathrm{CF}_t}
         {(1+r)^{\Delta_t/365}},
\qquad
\Delta_t = \operatorname{Tage}(d_0,d_t)
\]

definiert. \(d_0\) bezeichnet den Bewertungs- und Investitionszeitpunkt;
\(d_t\) ist das Zahlungsdatum des Cashflows der Periode \(t\). Bei
festgelegtem Diskontsatz stellt \(\operatorname{XNPV}(r)\) die explizit
auswertbare Barwertgröße des Modells dar.

\hypertarget{interner-zinsfuuxdf-als-implizite-renditekennzahl}{%
\subsubsection{2.2.2 Interner Zinsfuß als implizite
Renditekennzahl}\label{interner-zinsfuuxdf-als-implizite-renditekennzahl}}

Der auf datierten Zahlungszeitpunkten basierende interne Zinsfuß (XIRR)
ist nicht durch eine eigenständige Summenformel definiert. Er ergibt
sich als eine Nullstelle der zuvor definierten Barwertfunktion:

\[
r^{*} \in
\left\{r>-1\;\middle|\;\operatorname{XNPV}(r)=0\right\}.
\]

Damit beruhen Nettobarwert und interner Zinsfuß auf derselben
Cashflow-Zeitreihe. Der Nettobarwert bewertet diese Zeitreihe bei einem
exogen vorgegebenen Renditeanspruch; der interne Zinsfuß bestimmt
demgegenüber den endogenen Diskontsatz, bei dem ihr Barwert null
beträgt.

\hypertarget{periodische-equity-cashflows}{%
\subsubsection{2.2.3 Periodische
Equity-Cashflows}\label{periodische-equity-cashflows}}

Am Investitionszeitpunkt gilt mit Investitionsvolumen \(I\),
Fremdkapital \(D\) und Eigenkapitalquote \(e\):

\[
\mathrm{CF}_0 = -I + D = -eI,
\qquad
D=(1-e)I.
\]

Für die Betriebsperioden ergibt sich der Equity-Cashflow aus Erlösen,
Betriebskosten, Zinsen, Ertragsteuern und Tilgung:

\[
\mathrm{CF}_t
  = R_t-C_t-Z_t-S_t-T_t,
\qquad t\geq 1.
\]

Die Barwertfunktion kann damit zunächst in der folgenden aggregierten
Form geschrieben werden:

\[
\boxed{
\operatorname{XNPV}(r)
  = -eI
    + \sum_{t=1}^{N}
      \frac{R_t-C_t-Z_t-S_t-T_t}
           {(1+r)^{\Delta_t/365}}
}
\]

\hypertarget{substitution-von-erluxf6s-und-energieertrag}{%
\subsubsection{2.2.4 Substitution von Erlös und
Energieertrag}\label{substitution-von-erluxf6s-und-energieertrag}}

Für die Behandlung negativer Strompreise wird der Einspeisefaktor
\(q_t\) definiert als

\[
q_t =
\begin{cases}
1, & \text{im Modus Marktwert},\\
1-\nu_t, & \text{im Modus Abregelung}.
\end{cases}
\]

Die wirksame Marktprämie je Kilowattstunde beträgt

\[
p_t
  = \mathbf{1}_{[t\leq F]}\,(z-m_t)^+.
\]

Der Gesamterlös folgt damit aus

\[
R_t
  = \frac{E_t}{100}
    \left[q_t m_t+(1-\nu_t)p_t\right],
\]

wobei Energieertrag und nominaler Marktwert durch

\[
\begin{aligned}
E_t
  &= P h(1-d)^{t-1}\pi_t(1-\sigma),\\
m_t
  &= m_t^{\mathrm{real}}(1+\iota)^{y_t-y_B},\\
y_t
  &= y_0+t-1
\end{aligned}
\]

bestimmt werden. Nach dieser Substitution lautet die Barwertfunktion

\[
\begin{aligned}
\operatorname{XNPV}(r)
={}&-eI+\sum_{t=1}^{N}\frac{1}{(1+r)^{\Delta_t/365}}
\Bigg[
\frac{P h(1-d)^{t-1}\pi_t(1-\sigma)}{100}\\
&\quad\cdot\left(
q_t m_t+(1-\nu_t)\mathbf{1}_{[t\leq F]}(z-m_t)^+
\right)
-C_t-Z_t-S_t-T_t
\Bigg].
\end{aligned}
\]

Aus der Gleichung ist die Wirkung der technischen und marktlichen
Parameter auf den Projektbarwert unmittelbar ableitbar. Die übrigen
Terme werden im Folgenden als eigenständige Teilmodelle spezifiziert.

\hypertarget{substitution-der-betriebskosten}{%
\subsubsection{2.2.5 Substitution der
Betriebskosten}\label{substitution-der-betriebskosten}}

Die gesamten Betriebskosten setzen sich aus leistungsbezogenen
Standardpositionen, Gemeindeabgabe, Direktvermarktung und Pacht
zusammen:

\[
\begin{aligned}
C_t
={}&\sum_j
\mathbf{1}_{[t\geq t_j^{\mathrm{start}}]}
 w_jP(1+g_j)^{(t-t_j^{\mathrm{idx}})^+}\\
&+E_t c_{\mathrm{gem}}\Theta_t
 +C_t^{\mathrm{dv}}
 +C_t^{\mathrm{pacht}},\\
\Theta_t
={}&(1+\kappa)^{(t-1)^+}.
\end{aligned}
\]

Je nach vertraglicher Parametrisierung können einzelne Kostenpositionen
von der installierten Leistung, der erzeugten Energiemenge, dem
Marktwert oder dem Umsatz abhängen. Die vollständigen
Fallunterscheidungen werden in Kapitel 8 definiert.

\hypertarget{substitution-von-finanzierung-und-ertragsteuern}{%
\subsubsection{2.2.6 Substitution von Finanzierung und
Ertragsteuern}\label{substitution-von-finanzierung-und-ertragsteuern}}

Zinsaufwand und Darlehensstand ergeben sich aus

\[
Z_t=B_t i f_t,
\qquad
B_{t+1}=(B_t-T_t)^+.
\]

Die Steuerzahlung ist eine Funktion des Ergebnisses nach Betriebskosten
und Zinsen, der Abschreibung, des Freibetrags und des nutzbaren
Verlustvortrags:

\[
\begin{aligned}
G_t
  &=R_t-C_t-Z_t-A_t-\Phi,\\
U_t
  &=\min(V_t,\gamma G_t)\,
    \mathbf{1}_{[G_t>0]},\\
S_t
  &=\tau(G_t-U_t)^+,\\
V_{t+1}
  &=V_t-U_t+(-G_t)^+.
\end{aligned}
\]

Damit ist die Bewertungsfunktion vollständig auf Eingabeparameter,
Szenariokurven und periodenabhängige Zustandsgrößen zurückgeführt. Die
Kapitel 4 bis 10 spezifizieren die einzelnen Teilmodelle in der
Reihenfolge ihrer Implementierung; Kapitel 11 führt die Ergebnisse zum
Equity-Cashflow zusammen, Kapitel 12 wertet die Barwert- und
Renditekennzahlen aus.

\hypertarget{reihenfolge-und-fachliche-abhuxe4ngigkeiten}{%
\subsection{2.3 Reihenfolge und fachliche
Abhängigkeiten}\label{reihenfolge-und-fachliche-abhuxe4ngigkeiten}}

Die Reihenfolge der Rechenmodule folgt aus den fachlichen Abhängigkeiten
der jeweiligen Ausgangsgrößen:

\begin{longtable}[]{@{}
  >{\raggedright\arraybackslash}p{(\columnwidth - 6\tabcolsep) * \real{0.2500}}
  >{\raggedright\arraybackslash}p{(\columnwidth - 6\tabcolsep) * \real{0.2500}}
  >{\raggedright\arraybackslash}p{(\columnwidth - 6\tabcolsep) * \real{0.2500}}
  >{\raggedright\arraybackslash}p{(\columnwidth - 6\tabcolsep) * \real{0.2500}}@{}}
\toprule\noalign{}
\begin{minipage}[b]{\linewidth}\raggedright
Schritt
\end{minipage} & \begin{minipage}[b]{\linewidth}\raggedright
Modul
\end{minipage} & \begin{minipage}[b]{\linewidth}\raggedright
Erforderliche Eingangsgrößen
\end{minipage} & \begin{minipage}[b]{\linewidth}\raggedright
Fachliche Abhängigkeit
\end{minipage} \\
\midrule\noalign{}
\endhead
\bottomrule\noalign{}
\endlastfoot
1 & \texttt{timeline} & Inbetriebnahmedatum, Betriebsdauer & Definition
der Perioden und zeitanteiligen Faktoren \\
2 & \texttt{energy} & Zeitachse & Wirkung von Degradation und
unterjähriger Inbetriebnahme auf die Energiemenge \\
3 & \texttt{revenue} & Energiemenge, Kalenderjahr & Verknüpfung der
Produktion mit kalenderjahrbezogenen Marktpreiskurven \\
4 & \texttt{opex} & Energiemenge, Marktwert, Erlös & Berücksichtigung
mengen-, preis- und umsatzabhängiger Kostenpositionen \\
5 & \texttt{financing} & Investitions- und Kreditparameter & Ermittlung
des Schuldendienstes unabhängig vom operativen Ergebnis \\
6 & \texttt{tax} & Erlös, Betriebskosten, Zinsen, Investition &
Ableitung der steuerlichen Bemessungsgrundlage und Fortschreibung des
Verlustvortrags \\
7 & \texttt{cashflow} & sämtliche vorgelagerten Zeitreihen & Aggregation
zum Equity-Cashflow und zum DSCR \\
8 & \texttt{kpis} & datierte Cashflow-Zeitreihe & Berechnung der
Barwert-, Rendite- und Amortisationskennzahlen \\
\end{longtable}

Die Betriebskostenberechnung ist nach der Erlösberechnung angeordnet, da
Direktvermarktungskosten und Pacht in bestimmten Parametrisierungen vom
Marktwert beziehungsweise vom Umsatz abhängen. Eine umgekehrte
Abhängigkeit der Erlöse von den Betriebskosten ist nicht vorgesehen;
algebraische Zirkelbezüge werden dadurch ausgeschlossen.

\hypertarget{parametrisierungsebenen}{%
\subsection{2.4 Parametrisierungsebenen}\label{parametrisierungsebenen}}

Das Modell unterscheidet zwischen projektbezogenen und globalen
Parametern:

\begin{itemize}
\tightlist
\item
  Die \textbf{Projektparameter} (\texttt{PVProject}) enthalten
  projektspezifische Größen, insbesondere Leistung, spezifischen Ertrag,
  Investitionskosten, Standortkosten, Finanzierungskonditionen und
  Zuschlagswert.
\item
  Die \textbf{globalen Annahmen} (\texttt{GlobalAssumptions}) enthalten
  übergreifende Modellparameter, darunter Preiszeitreihen,
  Standardbetriebskosten, Förder- und Betriebsdauer, Steuerregeln,
  Degradation und Marktsystematik.
\end{itemize}

Die Funktion \texttt{resolve\_assumptions()} führt beide Ebenen zum
vollständig spezifizierten Parametersatz \texttt{EffectiveAssumptions}
zusammen. Sämtliche nachgelagerten Rechenmodule verwenden ausschließlich
diesen aufgelösten Satz. Dadurch bleibt jede Ergebnisgröße eindeutig auf
ihre Parametrisierung zurückführbar.

\hypertarget{marktsystemabhuxe4ngige-regelkonfiguration}{%
\subsection{2.5 Marktsystemabhängige
Regelkonfiguration}\label{marktsystemabhuxe4ngige-regelkonfiguration}}

Der Parameter \texttt{markt\_system} konfiguriert mehrere
länderspezifische Regeln als konsistentes Ausgangspaket:

\begin{longtable}[]{@{}
  >{\raggedright\arraybackslash}p{(\columnwidth - 4\tabcolsep) * \real{0.3333}}
  >{\raggedright\arraybackslash}p{(\columnwidth - 4\tabcolsep) * \real{0.3333}}
  >{\raggedright\arraybackslash}p{(\columnwidth - 4\tabcolsep) * \real{0.3333}}@{}}
\toprule\noalign{}
\begin{minipage}[b]{\linewidth}\raggedright
Regel
\end{minipage} & \begin{minipage}[b]{\linewidth}\raggedright
Österreich (EAG)
\end{minipage} & \begin{minipage}[b]{\linewidth}\raggedright
Deutschland (EEG)
\end{minipage} \\
\midrule\noalign{}
\endhead
\bottomrule\noalign{}
\endlastfoot
Entfall der Prämie bei negativen Preisen & ab 6 zusammenhängenden
Stunden & ab 1 Stunde \\
Ertragsbesteuerung & Körperschaftsteuer mit AfA, Freibetrag und
Verlustvortrag & Gewerbesteuer mit Freibetrag, ohne Verlustvortrag im
Referenzmodell \\
Zinsmethode im Anlaufjahr & taggenau Act/365 & kaufmännisch 30/360 \\
Herkunft des anzulegenden Wertes & empirisches Ausschreibungsmodell nach
Kapitel 15 & manuelle Vorgabe des erwarteten Zuschlags \\
\end{longtable}

Die Einzelparameter bleiben nach Auswahl des Marktsystems veränderbar.
Der Länderschalter stellt somit eine konsistente Vorbelegung dar, ohne
die Parametrisierung der Einzelregeln einzuschränken. Die grundlegende
Cashflow-Systematik bleibt in beiden Marktsystemen unverändert.

\hypertarget{notation-einheiten-und-konventionen}{%
\section{3 Notation, Einheiten und
Konventionen}\label{notation-einheiten-und-konventionen}}

\hypertarget{zeitindex}{%
\subsection{3.1 Zeitindex}\label{zeitindex}}

Der Index \(t\) bezeichnet das \textbf{Betriebsjahr}, gezählt ab der
Inbetriebnahme:

\[ t \in \{0, 1, 2, \dots, N\}, \qquad N = \text{Betriebsdauer in Jahren} \]

\begin{itemize}
\tightlist
\item
  \(t = 0\) ist der \textbf{Investitionszeitpunkt}. In diesem Jahr gibt
  es keine Produktion, keine Erlöse und keine Kosten, sondern
  ausschließlich den Investitionsabfluss und die Kreditaufnahme.
\item
  \(t \geq 1\) sind die \textbf{Betriebsjahre}. Betriebsjahr \(t\) endet
  stets am

  \begin{enumerate}
  \def\labelenumi{\arabic{enumi}.}
  \setcounter{enumi}{30}
  \tightlist
  \item
    Dezember des Kalenderjahres \(y_t\).
  \end{enumerate}
\end{itemize}

Das zugehörige \textbf{Kalenderjahr} ergibt sich aus dem
Inbetriebnahmejahr \(y_0\):

\[ y_t = y_0 + t - 1 \]

Diese Unterscheidung ist zentral: Preiskurven sind nach Kalenderjahr
indiziert, Degradation und Kostenindexierung nach Betriebsjahr.

\hypertarget{symbolverzeichnis}{%
\subsection{3.2 Symbolverzeichnis}\label{symbolverzeichnis}}

\begin{longtable}[]{@{}
  >{\raggedright\arraybackslash}p{(\columnwidth - 6\tabcolsep) * \real{0.2500}}
  >{\raggedright\arraybackslash}p{(\columnwidth - 6\tabcolsep) * \real{0.2500}}
  >{\raggedright\arraybackslash}p{(\columnwidth - 6\tabcolsep) * \real{0.2500}}
  >{\raggedright\arraybackslash}p{(\columnwidth - 6\tabcolsep) * \real{0.2500}}@{}}
\toprule\noalign{}
\begin{minipage}[b]{\linewidth}\raggedright
Symbol
\end{minipage} & \begin{minipage}[b]{\linewidth}\raggedright
Bedeutung
\end{minipage} & \begin{minipage}[b]{\linewidth}\raggedright
Einheit
\end{minipage} & \begin{minipage}[b]{\linewidth}\raggedright
Quelle
\end{minipage} \\
\midrule\noalign{}
\endhead
\bottomrule\noalign{}
\endlastfoot
\(P\) & Nennleistung (installiert, DC) & kWp & Projekt \\
\(h\) & Spezifischer Ertrag (Vollbenutzungsstunden) & kWh/kWp &
Projekt \\
\(d\) & Degradation je Jahr & -- & global \\
\(\sigma\) & Sicherheitsabschlag auf die Menge & -- & global \\
\(\pi_t\) & Anteilsfaktor der Periode \(t\) & -- & Zeitachse \\
\(E_t\) & Stromproduktion im Jahr \(t\) & kWh & Schritt 2 \\
\(m^{\mathrm{real}}_t\) & Marktwert Solar, real & ct/kWh &
Szenariokurve \\
\(m_t\) & Marktwert Solar, nominal & ct/kWh & Schritt 3 \\
\(b_t\) & Großhandelspreis (Baseload), nominal & ct/kWh &
Szenariokurve \\
\(z\) & EAG-Zuschlagswert, effektiv & ct/kWh & Projekt + Regel \\
\(F\) & Förderdauer & Jahre & global \\
\(s_t\) & Vergütungssatz & ct/kWh & Schritt 3 \\
\(\nu_t\) & Erzeugungsanteil in Stunden negativer Preise & -- &
Szenariokurve \\
\(w\) & Gewichtung des Negativstunden-Effekts & -- & global \\
\(R_t\) & Erlös gesamt & € & Schritt 3 \\
\(\theta_j\) & Einspeisekurve: Anteil des Monats \(j\) an der
Jahreserzeugung & -- & global \\
\(E_{t,j}\) & Stromproduktion im Jahr \(t\), Monat \(j\) & kWh & Schritt
2 \\
\(u_t\) & Rückzahlung je kWh (Überförderung) & ct/kWh & Schritt 3 \\
\(\beta\) & Toleranzband über dem anzulegenden Wert & -- & global \\
\(\rho\) & Rückzahlungsanteil oberhalb des Bandes & -- & global \\
\(P_{\min}\) & Engpassleistung, ab der zurückzuzahlen ist & MW &
global \\
\(\alpha\) & Anteil der Menge unter PPA & -- & Projekt \\
\(q_0\), \(q_t\) & PPA-Preis (vereinbart / im Jahr \(t\)) & €/MWh,
ct/kWh & Projekt \\
\(t_{\mathrm{PPA}}\), \(L\) & Startjahr und Laufzeit des PPA & Jahre &
Projekt \\
\(\eta\) & Indexierung des PPA-Preises & -- & Projekt \\
\(C_t\) & Betriebskosten gesamt & € & Schritt 4 \\
\(\kappa\) & Allgemeine Kosteninflation & -- & global \\
\(\iota\) & Inflation der Marktpreiskurven & -- & global \\
\(y_B\) & Preisbasisjahr der Marktpreiskurven & -- & global \\
\(I\) & Investitionsvolumen (CAPEX) & € & Projekt \\
\(e\) & Eigenkapitalquote & -- & Projekt \\
\(D\) & Kreditsumme (Fremdkapital) & € & Schritt 5 \\
\(i\) & Fremdkapitalzins & -- & Projekt \\
\(n\) & Kreditlaufzeit (Anzahl Tilgungsraten) & Jahre & global \\
\(Z_t\) & Zinsaufwand & € & Schritt 5 \\
\(T_t\) & Tilgung & € & Schritt 5 \\
\(B_t\) & Darlehensstand zu Jahresbeginn & € & Schritt 5 \\
\(A_t\) & Abschreibung (AfA) & € & Schritt 6 \\
\(V_t\) & Verlustvortragsbestand zu Jahresbeginn & € & Schritt 6 \\
\(\gamma\) & Verrechnungsgrenze des Verlustvortrags & -- & global \\
\(\tau\) & Effektiver Ertragsteuersatz & -- & global \\
\(S_t\) & Steuerzahlung & € & Schritt 6 \\
\(\mathrm{CF}_t\) & Equity-Cashflow gesamt & € & Schritt 7 \\
\(r\) & Diskontsatz & -- & Einstellung \\
\(s_{\mathrm{trap}}\) & Cash-Trap-Schwelle des DSCR & -- & global \\
\(s_{\mathrm{eod}}\) & Event-of-Default-Schwelle des DSCR & -- &
global \\
\(N_t\) & Erforderlicher Eigenkapitalnachschuss & € & Schritt 7 \\
\end{longtable}

\hypertarget{einheiten-und-umrechnungen}{%
\subsection{3.3 Einheiten und
Umrechnungen}\label{einheiten-und-umrechnungen}}

Das Modell rechnet \textbf{intern konsequent in kWh und ct/kWh} und
wandelt an den Rändern:

\[ \text{Erlös in €} = \frac{\text{Menge in kWh} \cdot \text{Satz in ct/kWh}}{100} \]

Eingaben, die üblicherweise je MWh erfasst werden (Gemeindeabgabe,
Direktvermarktungskosten), werden bereits bei der Parameterauflösung
umgerechnet:

\[ c^{\mathrm{kWh}} = \frac{c^{\mathrm{MWh}}}{1000} \]

Prozentangaben sind im Modell durchgängig \textbf{Brüche} in \([0,1]\)
(0,02 = 2 \%). Eine Ausnahme bildet der Gewerbesteuer-Hebesatz, dessen
natürlicher Wertebereich (200 bis 900) nicht in eine 0-bis-1-Konvention
passt; er wird als Prozentzahl geführt und in der Steuerformel durch 100
geteilt.

\hypertarget{vorzeichenkonvention}{%
\subsection{3.4 Vorzeichenkonvention}\label{vorzeichenkonvention}}

\begin{longtable}[]{@{}
  >{\raggedright\arraybackslash}p{(\columnwidth - 4\tabcolsep) * \real{0.3333}}
  >{\raggedright\arraybackslash}p{(\columnwidth - 4\tabcolsep) * \real{0.3333}}
  >{\raggedright\arraybackslash}p{(\columnwidth - 4\tabcolsep) * \real{0.3333}}@{}}
\toprule\noalign{}
\begin{minipage}[b]{\linewidth}\raggedright
Größe
\end{minipage} & \begin{minipage}[b]{\linewidth}\raggedright
Vorzeichen
\end{minipage} & \begin{minipage}[b]{\linewidth}\raggedright
Bemerkung
\end{minipage} \\
\midrule\noalign{}
\endhead
\bottomrule\noalign{}
\endlastfoot
Erlöse & positiv & \\
Betriebskosten, Zinsen, Steuern & positiv erfasst, negativ verrechnet &
die Zeitreihe führt Beträge, die Cashflow-Formel zieht sie ab \\
Investition \(\mathrm{CF}^{\mathrm{inv}}_0\) & negativ & Abfluss \\
Kreditaufnahme \(\mathrm{CF}^{\mathrm{fin}}_0\) & positiv & Zufluss \\
Tilgung \(\mathrm{CF}^{\mathrm{fin}}_t,\ t \geq 1\) & negativ &
Abfluss \\
\end{longtable}

\hypertarget{kappungen-und-randwertbehandlung}{%
\subsection{3.5 Kappungen und
Randwertbehandlung}\label{kappungen-und-randwertbehandlung}}

Drei Operationen treten wiederholt auf und sind einheitlich definiert:

\[ \mathrm{clip}(x, a, b) = \min\left(\max(x, a),\, b\right) \]

\[ (x)^{+} = \max(x, 0) \]

\[ \mathbf{1}_{[\text{Bedingung}]} = 1, \text{ wenn die Bedingung gilt, sonst } 0 \]

Für Kurvenzugriffe außerhalb des definierten Bereichs gilt durchgängig
\textbf{Klemmen statt Extrapolieren} (Abschnitt 7.2).

\hypertarget{schritt-0-aufluxf6sung-der-parameter}{%
\section{4 Schritt 0 -- Auflösung der
Parameter}\label{schritt-0-aufluxf6sung-der-parameter}}

\textbf{Zweck.} Aus Projektmaske und globalen Annahmen entsteht ein
vollständiger, widerspruchsfreier Parametersatz.

\textbf{Codestelle.} \nolinkurl{engine/pipeline.py},
\texttt{resolve\_assumptions()}.

\hypertarget{uxfcbernahme-umrechnung-geschuxe4ftsregeln}{%
\subsection{4.1 Übernahme, Umrechnung,
Geschäftsregeln}\label{uxfcbernahme-umrechnung-geschuxe4ftsregeln}}

Der Merge ist überwiegend eine 1:1-Übernahme. Vier Stellen enthalten
echte Rechenlogik:

\textbf{(a) Geschäftsregel Konventionell-Abschlag.} Konventionelle
Freiflächenanlagen erhalten gegenüber Agri-PV einen Abschlag von 25 \%
auf den EAG-Zuschlagswert:

\[ z = z_{\mathrm{Gebot}} \cdot \left(1 - \delta_{\mathrm{konv}} \cdot \mathbf{1}_{[\text{Anlagentyp} = \text{konventionell}]}\right), \qquad \delta_{\mathrm{konv}} = 0{,}25 \]

Der Abschlag ist als benannte Konstante
(\texttt{KONVENTIONELL\_ZUSCHLAG\_ABSCHLAG\_PCT}) implementiert und
nicht als Eingabeparameter veränderbar. Damit wird er als Geschäftsregel
und nicht als projektspezifische Annahme behandelt.

\textbf{(b) Einheitenumrechnung} der produktionsbasierten Sätze:

\[ c_{\mathrm{gem}} = \frac{c^{\mathrm{MWh}}_{\mathrm{gem}}}{1000}, \qquad c_{\mathrm{dv}} = \frac{c^{\mathrm{MWh}}_{\mathrm{dv}}}{1000} \]

\textbf{(c) Auswahl der Negativmengen-Kurve.} Jedes Marktpreisszenario
führt zwei Zeitreihen des Erzeugungsanteils in Stunden negativer Preise
-- eine für die 6-Stunden-Regel, eine für die 1-Stunden-Regel. Die
global gewählte Regel entscheidet, welche in den Parametersatz
übernommen wird. Da die 1-Stunden-Regel einen größeren Stundenumfang
erfasst, weist die zugehörige Mengenkurve mindestens gleich hohe Werte
auf.

\textbf{(d) Szenario-Rückfall.} Existiert das im Projekt hinterlegte
Szenario nicht mehr, wird das erste verfügbare Szenario verwendet;
existiert gar keines, ein leeres Szenario mit Marktwert 0. Dadurch
bleibt die Bewertung auch nach Umbenennung oder Löschung eines
referenzierten Szenarios ausführbar. Der verwendete Rückfallwert ist im
aufgelösten Parametersatz transparent ausgewiesen.

\hypertarget{investitionsvolumen}{%
\subsection{4.2 Investitionsvolumen}\label{investitionsvolumen}}

Das Investitionsvolumen ist die Summe der neun festen Kostenkategorien
der Projektmaske und beliebig vieler frei benannter Zusatzpositionen:

\[ I = \sum_{c\, \in\, \mathcal{C}} I_c + \sum_{z\, \in\, \mathcal{Z}} I_z \]

\[ \mathcal{C} = \{\text{EPC},\ \text{Netzanschluss},\ \text{Trasse},\ \text{Widmung},\ \text{Genehmigung},\ \text{Sonstige externe},\ \text{AGM},\ \text{M+A},\ \text{Pönalepuffer}\} \]

\(\mathcal{Z}\) ist die je Projekt frei definierbare Menge zusätzlicher
Investitionspositionen. Jede Position besteht aus einer Bezeichnung und
einem Betrag; die Bezeichnung muss nichtleer sein und darf keine der
reservierten Spaltenbezeichnungen der Cashflow-Zeitreihe tragen, damit
die Ausgabestrukturen eindeutig bleiben. Rechnerisch sind die
Zusatzpositionen den festen Kategorien gleichgestellt -- sie erhöhen
ausschließlich \(I\) und wirken damit über die Finanzierungs-,
Abschreibungs- und Cashflow-Kette identisch. Ist \(\mathcal{Z}\) leer,
reduziert sich die Formel auf die neun Standardkategorien.

Alle Kategorien werden als \textbf{Gesamtbeträge in Euro} erfasst; die
Projektmaske erlaubt je Feld wahlweise die Eingabe als spezifischer Wert
(€/kWp), der intern unmittelbar mit \(P\) multipliziert wird. Die
spezifische Größe \(I/P\) wird nur zur Anzeige gebildet.

\hypertarget{vollstuxe4ndigkeitspruxfcfung}{%
\subsection{4.3
Vollständigkeitsprüfung}\label{vollstuxe4ndigkeitspruxfcfung}}

Zwei Konsistenzbedingungen werden bereits bei der Validierung erzwungen:

\begin{itemize}
\tightlist
\item
  Im Steuermodus mit AfA muss eine Nutzungsdauer gesetzt sein (sonst
  wäre \(A_t\) undefiniert).
\item
  Alle Anteilsgrößen liegen in \([0,1]\), Leistung und spezifischer
  Ertrag sind strikt positiv.
\end{itemize}

\textbf{Ausgang.} \texttt{EffectiveAssumptions} mit rund 40 Feldern als
einheitliche Eingangsstruktur sämtlicher nachgelagerter Rechenmodule.

\hypertarget{schritt-1-zeitachse}{%
\section{5 Schritt 1 -- Zeitachse}\label{schritt-1-zeitachse}}

\textbf{Zweck.} Festlegen der Perioden und ihrer Anteilsfaktoren.

\textbf{Codestelle.} \nolinkurl{engine/timeline.py},
\texttt{build\_timeline()} und \texttt{erstjahr\_zins\_pro\_rata()}.

\hypertarget{perioden}{%
\subsection{5.1 Perioden}\label{perioden}}

Betriebsjahr \(t\) läuft von \(\mathrm{start}_t\) bis
\(\mathrm{ende}_t\) mit

\[ \mathrm{ende}_t = 31.12.(y_0 + t - 1), \qquad \mathrm{start}_1 = \text{Inbetriebnahmedatum}, \qquad \mathrm{start}_{t+1} = 1.1.(y_0 + t) \]

Jede Periode außer der ersten ist damit ein volles Kalenderjahr. Die
erste Periode reicht vom Inbetriebnahmedatum bis zum Jahresende und ist
bei unterjähriger Inbetriebnahme kürzer.

\hypertarget{anteilsfaktor-der-produktion}{%
\subsection{5.2 Anteilsfaktor der
Produktion}\label{anteilsfaktor-der-produktion}}

\[ \pi_t = \min\left(\frac{\mathrm{Tage}(\mathrm{start}_t,\ \mathrm{ende}_t) + 1}{365},\ 1\right) \]

Der Faktor ist für alle vollen Kalenderjahre gleich 1 (auch in
Schaltjahren, wegen der Kappung bei 1) und nur im Anlaufjahr kleiner. Er
skaliert die Produktionsmenge des ersten Jahres und damit indirekt alle
mengenabhängigen Erlöse und Kosten.

\begin{quote}
\textbf{Annahme.} Betriebsperioden sind Kalenderjahre. Der Sonderfall
„Vertragsende am Jahrestag der Inbetriebnahme statt am Jahresende`` ist
nicht abgebildet (siehe Kapitel 17).
\end{quote}

\hypertarget{anteilsfaktor-der-zinsen}{%
\subsection{5.3 Anteilsfaktor der
Zinsen}\label{anteilsfaktor-der-zinsen}}

Für den Zinsaufwand des Anlaufjahres ist die Zinsmethode wählbar. Sie
wirkt sich ausschließlich bei unterjähriger Inbetriebnahme aus; bei
Inbetriebnahme am 1. Januar liefern beide Methoden exakt 1.

\textbf{Österreich, taggenau (act/365):} identisch zum
Produktionsfaktor,

\[ f^{\mathrm{act/365}} = \min\left(\frac{\mathrm{Tage}(\text{IBN},\ 31.12.y_0) + 1}{365},\ 1\right) \]

\textbf{Deutsche kaufmännische Methode (30/360):} jeder Restmonat des
Anlaufjahres einschließlich des Inbetriebnahmemonats zählt pauschal mit
30 Tagen, das Jahr mit 360:

\[ f^{30/360} = \frac{(13 - M) \cdot 30}{360} = \frac{13 - M}{12}, \qquad M = \text{Inbetriebnahmemonat} \]

Beispiel: Inbetriebnahme im September (\(M = 9\)) ergibt
\(f^{30/360} = 4/12 = 0{,}3\overline{3}\), während act/365 vom 1.9. bis
31.12. mit \(122/365 = 0{,}334\) rechnet -- nahe beieinander, aber nicht
identisch.

\textbf{Ausgang.} Zeitachse mit \(t\), Periodengrenzen, \(\pi_t\) und
der Markierung des letzten Jahres.

\hypertarget{schritt-2-energieertrag}{%
\section{6 Schritt 2 -- Energieertrag}\label{schritt-2-energieertrag}}

\textbf{Zweck.} Die eingespeiste Strommenge je Betriebsjahr.

\textbf{Codestelle.} \nolinkurl{engine/energy.py},
\texttt{calculate\_energy\_production()}.

\hypertarget{vorschrift}{%
\subsection{6.1 Vorschrift}\label{vorschrift}}

\[ E^{\mathrm{basis}} = P \cdot h \]

\[ \phi_t = (1 - d)^{\,t-1} \]

\[ E_t = E^{\mathrm{basis}} \cdot \phi_t \cdot \pi_t \cdot (1 - \sigma) \]

\hypertarget{erluxe4uterung}{%
\subsection{6.2 Erläuterung}\label{erluxe4uterung}}

\begin{itemize}
\tightlist
\item
  \textbf{Basisproduktion} \(E^{\mathrm{basis}}\) ist der Ertrag eines
  vollen ersten Betriebsjahres ohne Degradation und ohne Abschlag; sie
  ist das Produkt aus installierter Leistung und spezifischem Ertrag.
\item
  \textbf{Degradationsfaktor} \(\phi_t\) ist geometrisch mit Exponent
  \(t-1\): Das erste Betriebsjahr ist definitionsgemäß degradationsfrei
  (\(\phi_1 = 1\)), der Modulalterungseffekt wirkt ab dem zweiten Jahr.
  Bei \(d = 0{,}25\,\%\) und \(t = 30\) ergibt sich
  \(\phi_{30} = 0{,}9975^{29} = 0{,}9298\), entsprechend einem
  Mengenrückgang von rund 7 \% bis zum Ende der Betriebsdauer.
\item
  \textbf{Anteilsfaktor} \(\pi_t\) kürzt das Anlaufjahr (Abschnitt 5.2).
\item
  \textbf{Sicherheitsabschlag} \(\sigma\) ist ein pauschaler Abschlag
  auf die Ertragsprognose (z. B. für P90-Betrachtungen). Er wirkt
  multiplikativ auf alle Jahre gleich.
\end{itemize}

\begin{quote}
Der Sicherheitsabschlag wird \textbf{nach} Degradation und Anteilsfaktor
angewendet. Da alle drei Faktoren multiplikativ sind, ist die
Reihenfolge rechnerisch ohne Wirkung; die Trennung hält die Zeitreihe
aber interpretierbar (der ausgewiesene Degradationsfaktor enthält keinen
Risikoabschlag).
\end{quote}

\hypertarget{monatsaufteilung-optionale-zeitaufluxf6sung}{%
\subsection{6.3 Monatsaufteilung (optionale
Zeitauflösung)}\label{monatsaufteilung-optionale-zeitaufluxf6sung}}

Wird in den globalen Annahmen die \textbf{Monatsauflösung} gewählt, wird
die Jahresmenge über eine \textbf{Einspeisekurve}
\(\theta_1,\dots,\theta_{12}\) verteilt. Sie gibt den Anteil der
Jahreserzeugung je Kalendermonat an und wird vor der Anwendung normiert,
damit gerundete Prozentangaben die Jahresmenge nicht verändern:

\[ \hat{\theta}_j = \frac{\theta_j}{\sum_{i=1}^{12}\theta_i}, \qquad j = 1,\dots,12 \]

\[ E_{t,j} = E^{\mathrm{basis}} \cdot \phi_t \cdot (1-\sigma) \cdot \hat{\theta}_j \cdot \mathbf{1}_{[\,t>1 \;\vee\; j \geq j_0\,]} \]

Dabei ist \(j_0\) der Inbetriebnahmemonat. Der Anteilsfaktor \(\pi_t\)
des Anlaufjahres entfällt in dieser Auflösung: Statt eines taggenauen
Bruchteils zählen genau die Monate ab Inbetriebnahme mit ihrem
tatsächlichen Ertragsanteil. Für eine im Juli in Betrieb gehende Anlage
sind das nicht die halben, sondern rund 60 \% der Jahresmenge -- die
ertragreichen Monate liegen im Sommer. Für alle weiteren Jahre gilt

\[ \sum_{j=1}^{12} E_{t,j} = E_t, \]

die Monatsebene ist also eine reine Verteilung. Alle nachgelagerten
Schritte (Betriebskosten, Finanzierung, Steuern, Kennzahlen) rechnen
unverändert auf Jahresscheiben; die Monatsebene wirkt ausschließlich in
Schritt 3 (Abschnitt 7.9).

\textbf{Ausgang.} \(\phi_t\) und \(E_t\) (kWh) je Betriebsjahr, in der
Monatsauflösung zusätzlich \(E_{t,j}\).

\hypertarget{schritt-3-erluxf6se-nach-dem-marktpruxe4mienmodell}{%
\section{7 Schritt 3 -- Erlöse nach dem
Marktprämienmodell}\label{schritt-3-erluxf6se-nach-dem-marktpruxe4mienmodell}}

\textbf{Zweck.} Vergütungssatz und Erlös je Betriebsjahr, aufgeteilt in
Markterlös und Marktprämie.

\textbf{Codestelle.} \nolinkurl{engine/revenue.py},
\texttt{calculate\_revenue()}.

In diesem Modul werden die zentralen marktseitigen Berechnungsregeln
zusammengeführt. Die Berechnung gliedert sich in fünf Teilschritte:
Kalenderjahr-Zuordnung, Kurvenzugriff, Inflationierung, Vergütungssatz,
Mengenwirkung negativer Preise.

Drei Erweiterungen setzen darauf auf und sind voneinander unabhängig:
das \textbf{Prämienmodell} (Abschnitt 7.7) bestimmt, was oberhalb des
anzulegenden Werts geschieht; die \textbf{hybride Vermarktung}
(Abschnitt 7.8) teilt die Menge zwischen PPA und Spotmarkt auf; die
\textbf{Monatsauflösung} (Abschnitt 7.9) führt Menge und Preis auf
Monatsebene zusammen. Die Abschnitte 7.1 bis 7.6 beschreiben den
Grundfall: einseitiger CfD, reine Merchant-Vermarktung, Jahresauflösung.

\hypertarget{kalenderjahr-zuordnung}{%
\subsection{7.1 Kalenderjahr-Zuordnung}\label{kalenderjahr-zuordnung}}

\[ y_t = y_0 + t - 1 \]

Alle Preis- und Mengenkurven der Szenarien sind nach \textbf{echtem
Kalenderjahr} indiziert (typisch 2025 bis 2060), nicht nach
Betriebsjahr. Zwei Projekte mit unterschiedlichem Inbetriebnahmejahr
greifen im selben Betriebsjahr folglich auf unterschiedliche
Kurvenpunkte zu. Dies entspricht der kalenderzeitbezogenen Struktur
einer Marktpreisprognose.

\hypertarget{kurvenzugriff-mit-klemmung}{%
\subsection{7.2 Kurvenzugriff mit
Klemmung}\label{kurvenzugriff-mit-klemmung}}

Für eine Kurve \(g\), die auf den Stützjahren
\(\{y_{\min}, \dots, y_{\max}\}\) definiert ist:

\[ g(y_t) = g\left(\mathrm{clip}(y_t,\ y_{\min},\ y_{\max})\right) \]

Liegt das Kalenderjahr außerhalb des Kurvenbereichs, wird der
nächstgelegene \textbf{Randwert} verwendet. Eine Extrapolation erfolgt
nicht, da eine lineare Fortschreibung des letzten Kurvensegments über
den beobachteten Zeitraum hinaus nicht durch die zugrunde liegende
Marktpreisstudie abgesichert wäre. Ist gar keine Kurve hinterlegt, gilt
\(g \equiv 0\).

\hypertarget{inflationierung-des-marktwertes}{%
\subsection{7.3 Inflationierung des
Marktwertes}\label{inflationierung-des-marktwertes}}

Marktwert-Solar-Kurven aus Marktpreisstudien sind \textbf{reale} Werte
auf der Preisbasis des Erscheinungsjahres \(y_B\). Für eine nominale
Cashflow-Rechnung werden sie inflationiert:

\[ m_t = m^{\mathrm{real}}_t \cdot (1 + \iota)^{\,y_t - y_B} \]

Für die Anwendung sind zwei Aspekte maßgeblich:

\begin{enumerate}
\def\labelenumi{\arabic{enumi}.}
\tightlist
\item
  Der Exponent verwendet das \textbf{tatsächliche} Kalenderjahr \(y_t\),
  nicht das eventuell geklemmte Nachschlagejahr. Wird jenseits des
  letzten Kurvenjahres mit dem letzten bekannten Realpreis
  fortgeschrieben, wird die allgemeine Preisniveauentwicklung weiterhin
  berücksichtigt.
\item
  Der \textbf{EAG-Zuschlagswert \(z\) bleibt nominal fix}. Er wird
  während der Förderdauer gesetzlich nicht indexiert. Der Vergleich
  zwischen Marktwert und Zuschlagswert findet damit korrekt zwischen
  zwei nominalen Größen statt.
\end{enumerate}

\hypertarget{verguxfctungssatz-gleitende-marktpruxe4mie}{%
\subsection{7.4 Vergütungssatz (gleitende
Marktprämie)}\label{verguxfctungssatz-gleitende-marktpruxe4mie}}

Während der Förderdauer \(F\) erhält der Betreiber den höheren der
beiden Werte, danach ausschließlich den Marktwert:

\[
\begin{aligned}
\widetilde{p}_t &= \left(z-m_t\right)^+,\\
p_t &= \mathbf{1}_{[t\leq F]}\,\widetilde{p}_t,\\
s_t &= m_t+p_t.
\end{aligned}
\]

Äquivalent und anschaulicher:

\[
s_t=
\begin{cases}
\max(m_t,z), & t\leq F,\\
m_t, & t>F.
\end{cases}
\]

\(\widetilde{p}_t\) ist die rechnerische Differenz zwischen
Zuschlagswert und Marktwert; \(p_t\) ist die \textbf{tatsächlich
wirksame Marktprämie je kWh}. Liegt der Marktwert unter dem
Zuschlagswert, wird die Differenz während der Förderdauer zugeschossen;
liegt er darüber, ist die Prämie null. Nach Ablauf der Förderdauer ist
\(p_t\) definitionsgemäß null.

\hypertarget{wirkung-negativer-strompreise}{%
\subsection{7.5 Wirkung negativer
Strompreise}\label{wirkung-negativer-strompreise}}

In Stunden negativer Preise entfällt die Marktprämie. Modelliert wird
das über den Anteil \(\nu^{\mathrm{roh}}_t\) der
\textbf{Erzeugungsmenge}, die in solche Stunden fällt (nicht über den
Anteil der Stunden selbst -- die Menge ist die wirtschaftlich relevante
Größe). Eine globale Gewichtung erlaubt das Ein- und Ausblenden des
Effekts für Vergleichsrechnungen:

\[ \nu_t = \nu^{\mathrm{roh}}_t \cdot w, \qquad w \in [0,1] \]

\(w = 0\) blendet den Effekt vollständig aus (volle Vergütung auch in
Stunden negativer Preise), \(w = 1\) entspricht der vollen gesetzlichen
Wirkung, wie sie in der Szenariokurve hinterlegt ist.

Für das Verhalten der Anlage in diesen Stunden gibt es zwei Modi. In
\textbf{beiden} entfällt die Prämie für den betroffenen Mengenanteil;
sie unterscheiden sich nur darin, ob der Markterlös erhalten bleibt.

\textbf{Modus „Marktwert`` (Anlage speist weiter ein):}

\[ R^{\mathrm{markt}}_t = \frac{E_t \cdot m_t}{100}, \qquad R^{\mathrm{pr\ddot{a}mie}}_t = \frac{E_t \cdot (1 - \nu_t) \cdot p_t}{100} \]

\textbf{Modus „Abregelung`` (Anlage wird abgeregelt):}

\[ R^{\mathrm{markt}}_t = \frac{E_t \cdot (1 - \nu_t) \cdot m_t}{100}, \qquad R^{\mathrm{pr\ddot{a}mie}}_t = \frac{E_t \cdot (1 - \nu_t) \cdot p_t}{100} \]

In beiden Fällen gilt

\[ R_t = R^{\mathrm{markt}}_t + R^{\mathrm{pr\ddot{a}mie}}_t \]

Nach Ablauf der Förderdauer ist \(p_t = 0\); der Modus „Marktwert`` hat
dann keine Wirkung mehr, während „Abregelung`` die Menge weiterhin
kürzt.

\hypertarget{interpretation-der-aufteilung}{%
\subsection{7.6 Interpretation der
Aufteilung}\label{interpretation-der-aufteilung}}

Die Zerlegung in \(R^{\mathrm{markt}}\) und \(R^{\mathrm{prämie}}\)
ermöglicht eine ökonomische Zuordnung der Erlöse zu Markt und Förderung.
Dadurch lässt sich insbesondere die Abhängigkeit des Projektwerts von
der Förderdauer quantifizieren. In der grafischen Darstellung entspricht
die Differenz zwischen Vergütungssatz und Marktwert der Marktprämie.

\hypertarget{marktpruxe4mienmodelle}{%
\subsection{7.7 Marktprämienmodelle}\label{marktpruxe4mienmodelle}}

Abschnitt 7.4 beschreibt den \textbf{einseitigen} Differenzvertrag: Der
Betreiber erhält den höheren von Marktwert und anzulegendem Wert. Zwei
weitere Vertragsformen sind wählbar; sie unterscheiden sich
ausschließlich im Bereich \(m_t > z\). Unterhalb des anzulegenden Werts
zahlen alle drei Modelle identisch auf.

Neben der Prämie \(p_t\) tritt dafür eine \textbf{Rückzahlung} \(u_t\)
je kWh -- eine Zahlung in die Gegenrichtung, an die Förderstelle:

\[
u_t = \mathbf{1}_{[t\leq F]} \cdot
\begin{cases}
0, & \text{einseitiger CfD},\\[2pt]
\left(m_t-z\right)^+, & \text{zweiseitiger CfD},\\[2pt]
\rho\cdot\left(m_t-z\,(1+\beta)\right)^+\cdot\mathbf{1}_{[P\geq P_{\min}]}, & \text{Toleranzband (EAG)}.
\end{cases}
\]

Der Vergütungssatz wird damit

\[ s_t = m_t + p_t - u_t. \]

\begin{itemize}
\tightlist
\item
  \textbf{Einseitiger CfD.} \(u_t \equiv 0\); der Betreiber behält jeden
  Übergewinn. Grundform der EAG-Marktprämie und Voreinstellung des
  Modells.
\item
  \textbf{Zweiseitiger CfD.} \(s_t = z\) während der gesamten
  Förderdauer: kein Preisrisiko nach unten, keine Preischance nach oben.
  Diese Form liegt der deutschen EEG-Reform 2027 zugrunde.
\item
  \textbf{Toleranzband (§ 10 EAG).} Eine Rückzahlung entsteht erst, wenn
  der Referenzmarktwert den anzulegenden Wert um mehr als \(\beta\)
  übersteigt; zurückzuzahlen ist der Anteil \(\rho\) des darüber
  liegenden Betrags. Die Pflicht gilt erst ab einer Engpassleistung
  \(P_{\min}\). Gesetzliche Werte für Photovoltaik: \(\beta = 40\,\%\),
  \(\rho = 66\,\%\), \(P_{\min} = 5\ \mathrm{MW}\). Alle drei sind
  einstellbar, da sie Gegenstand laufender Novellen sind; verglichen
  wird mit der Nennleistung \(P\) des Projekts.
\end{itemize}

Die Rückzahlung wird auf derselben Menge bemessen wie die Prämie
(Abschnitt 7.5) und mindert den Erlös:

\[ R^{\mathrm{r\ddot{u}ck}}_t = \frac{E_t\,(1-\nu_t)\,u_t}{100}, \qquad R_t = R^{\mathrm{markt}}_t + R^{\mathrm{pr\ddot{a}mie}}_t - R^{\mathrm{r\ddot{u}ck}}_t \]

Prämie und Rückzahlung werden getrennt ausgewiesen und nicht saldiert:
Es sind zwei Zahlungsrichtungen, und nur getrennt ist erkennbar, welcher
Teil des Ergebnisses aus einer Abschöpfung stammt.

\hypertarget{hybride-vermarktung-ppa-und-merchant}{%
\subsection{7.8 Hybride Vermarktung (PPA und
Merchant)}\label{hybride-vermarktung-ppa-und-merchant}}

Ein Anteil \(\alpha \in [0,1]\) der vermarkteten Menge wird zu einem
festen Preis \(q_t\) abgenommen, der Rest zum Marktwert verkauft. Der
PPA-Preis gilt zwischen dem Startjahr \(t_{\mathrm{PPA}}\) und dem Ende
der Laufzeit \(L\) und wird ab dem zweiten Vertragsjahr mit \(\eta\)
indexiert:

\[
q_t =
\begin{cases}
\dfrac{q_0}{10}\,(1+\eta)^{\,t-t_{\mathrm{PPA}}}, & t_{\mathrm{PPA}} \leq t \leq t_{\mathrm{PPA}}+L-1,\\[6pt]
0, & \text{sonst},
\end{cases}
\]

mit \(q_0\) in €/MWh und \(q_t\) in ct/kWh. Mit \(E^{\mathrm{verm}}_t\)
als der vermarkteten Menge nach Abschnitt 7.5 gilt

\[ R^{\mathrm{ppa}}_t = \frac{\alpha\,E^{\mathrm{verm}}_t\,q_t}{100}, \qquad R^{\mathrm{merch}}_t = \frac{(1-\alpha)\,E^{\mathrm{verm}}_t\,m_t}{100}, \]

\[ R^{\mathrm{markt}}_t = R^{\mathrm{ppa}}_t + R^{\mathrm{merch}}_t. \]

Entscheidend ist die Unabhängigkeit von der Förderung: Prämie \(p_t\)
und Rückzahlung \(u_t\) bemessen sich unverändert am
\textbf{Referenzmarktwert} \(m_t\), nicht am tatsächlich erzielten
Preis. Wer unterhalb des Referenzwerts verkauft, trägt die Differenz
selbst; wer darüber verkauft, behält sie. Das entspricht der
Konstruktion der gleitenden Marktprämie und ist zugleich der Grund,
warum ein PPA die Erlösverteilung verschiebt, ohne den Förderanspruch zu
berühren.

\hypertarget{monatsaufluxf6sung}{%
\subsection{7.9 Monatsauflösung}\label{monatsaufluxf6sung}}

In der Monatsauflösung (Abschnitt 6.3) tragen alle Kurven einen
Monatsindex: Marktwert \(m_{t,j}\) und Negativmengenanteil \(\nu_{t,j}\)
werden aus Monatsreihen des Szenarios gelesen. Liegt für ein
Kalenderjahr keine Monatsreihe vor, gilt dessen Jahreswert für alle
zwölf Monate -- ein Szenario bleibt damit auch dann rechenbar, wenn nur
ein Teil der Jahre in Monatsauflösung vorliegt.

Sämtliche Vorschriften der Abschnitte 7.3 bis 7.8 werden je
Monatsscheibe \((t,j)\) ausgewertet; die Inflationierung verwendet dabei
unverändert das Kalenderjahr \(y_t\). Die Verdichtung auf Jahreswerte
erfolgt anschließend:

\[ R_t = \sum_{j=1}^{12} R_{t,j}, \qquad
m_t = \frac{\sum_{j=1}^{12} E_{t,j}\, m_{t,j}}{\sum_{j=1}^{12} E_{t,j}} \]

Beträge werden \textbf{summiert}, Preise \textbf{mengengewichtet}
gemittelt. Die Gewichtung ist keine Feinheit: Der Jahresmarktwert einer
PV-Anlage ist der Wert, den \emph{ihre} Kilowattstunden erlösen. Ein
ungewichteter Monatsdurchschnitt wäre für Photovoltaik systematisch zu
hoch, weil die ertragsstarken Monate die preisschwachen sind.

Der wirtschaftliche Unterschied beider Auflösungen entsteht aus zwei
Effekten mit gegenläufigem Vorzeichen:

\begin{enumerate}
\def\labelenumi{\arabic{enumi}.}
\tightlist
\item
  \textbf{Erlösprofil.} Erzeugung und Preis sind negativ korreliert;
  eine Jahresrechnung überschätzt den Merchant-Erlös. Ist die
  Jahreskurve bereits ein erzeugungsgewichteter Marktwert Solar, ist
  dieser Effekt in ihr enthalten und die Monatsauflösung reproduziert
  ihn.
\item
  \textbf{Konvexität der Förderung.} Prämie und Rückzahlung sind
  abgeschnittene Funktionen von \(m_t\). Nach der jensenschen
  Ungleichung ist der Erwartungswert der Prämie über streuende
  Monatswerte größer als die Prämie des Mittelwerts -- der einseitige
  CfD gewinnt an Wert, je stärker die Monatspreise streuen. Erst die
  Monatsauflösung macht diesen Optionswert sichtbar.
\end{enumerate}

\textbf{Ausgang.} \(y_t\), \(m^{\mathrm{real}}_t\), \(m_t\), \(s_t\),
\(R_t\), \(R^{\mathrm{markt}}_t\), \(R^{\text{Prämie}}_t\),
\(R^{\mathrm{ppa}}_t\), \(R^{\mathrm{merch}}_t\),
\(R^{\mathrm{r\ddot{u}ck}}_t\) sowie -- sofern das Szenario ihn führt --
der nominale Großhandelspreis \(b_t\). Er geht in keine Erlösgröße ein;
er ist Bezugsgröße der Direktvermarktungskosten (Abschnitt 8.3) und im
Diagramm der Abstand, aus dem sich die Kannibalisierung ablesen lässt.

\hypertarget{schritt-4-betriebskosten}{%
\section{8 Schritt 4 -- Betriebskosten}\label{schritt-4-betriebskosten}}

\textbf{Zweck.} Alle laufenden Kosten je Betriebsjahr, sowohl als Summe
als auch aufgeschlüsselt nach Einzelposition.

\textbf{Codestelle.} \nolinkurl{engine/opex.py},
\texttt{calculate\_opex()}.

Die Betriebskosten setzen sich aus vier Gruppen zusammen, die
unterschiedlichen Bemessungsgrundlagen folgen:

\begin{longtable}[]{@{}
  >{\raggedright\arraybackslash}p{(\columnwidth - 4\tabcolsep) * \real{0.3333}}
  >{\raggedright\arraybackslash}p{(\columnwidth - 4\tabcolsep) * \real{0.3333}}
  >{\raggedright\arraybackslash}p{(\columnwidth - 4\tabcolsep) * \real{0.3333}}@{}}
\toprule\noalign{}
\begin{minipage}[b]{\linewidth}\raggedright
Gruppe
\end{minipage} & \begin{minipage}[b]{\linewidth}\raggedright
Bemessung
\end{minipage} & \begin{minipage}[b]{\linewidth}\raggedright
Indexierung
\end{minipage} \\
\midrule\noalign{}
\endhead
\bottomrule\noalign{}
\endlastfoot
Standardpositionen & €/kWp/Jahr & eigene, je Position sichtbare
Indexierung \\
Gemeindeabgabe & €/kWh & allgemeine Kosteninflation \\
Direktvermarktung & €/kWh oder \% vom Marktwert & Kosteninflation bzw.
implizit über den Marktwert \\
Pacht & €/kWp/Jahr oder \% vom Umsatz mit Mindestpacht &
Kosteninflation \\
\end{longtable}

\hypertarget{allgemeiner-inflationsfaktor}{%
\subsection{8.1 Allgemeiner
Inflationsfaktor}\label{allgemeiner-inflationsfaktor}}

Für alle Positionen ohne eigene Preislogik gilt ein gemeinsamer
Kosteninflationsfaktor. Eingaben verstehen sich als Preisstand bei
Inbetriebnahme, die Eskalation beginnt daher im zweiten Betriebsjahr:

\[ \Theta_t = (1 + \kappa)^{\,(t-1)^{+}} \]

\hypertarget{betriebskostenpositionen}{%
\subsection{8.2
Betriebskostenpositionen}\label{betriebskostenpositionen}}

Die Positionsliste eines Projekts entsteht aus zwei Quellen: den in den
Globalen Annahmen gepflegten Standardpositionen und den frei benannten
Zusatzpositionen des Projekts. Beide werden in
\texttt{resolve\_assumptions()} zu einer Liste verkettet --
Standardpositionen zuerst, danach die Zusatzpositionen. Die Reihenfolge
bestimmt zugleich die Stapelreihenfolge in der Kostendarstellung. Für
die Rechnung selbst besteht kein Unterschied: Jede Position durchläuft
dieselbe Vorschrift.

Jede Position \(j\) trägt einen spezifischen Basiswert \(w_j\)
(€/kWp/Jahr), ein Startjahr \(t^{\mathrm{start}}_j\), eine eigene
Indexrate \(g_j\) und ein Jahr, ab dem indexiert wird,
\(t^{\mathrm{idx}}_j\):

\[ C^{(j)}_t = \mathbf{1}_{[\,t\, \geq\, t^{\mathrm{start}}_j\,]} \cdot w_j \cdot P \cdot (1 + g_j)^{\left(t - t^{\mathrm{idx}}_j\right)^{+}} \]

Der Exponent ist bei \(0\) gekappt: Vor dem Indexierungsstartjahr gilt
der unveränderte Basiswert, es findet keine „Rückwärtsindexierung``
statt. Der Startzeitpunkt ist von der Indexierung getrennt, damit
Positionen abgebildet werden können, die erst später einsetzen (z. B.
eine Rücklage ab Jahr 11), ihren Preisstand aber ab Jahr 1
fortschreiben.

Positionen mit identischer Bezeichnung werden aggregiert. Dadurch bleibt
jede Bezeichnung in der Kostenaufschlüsselung eindeutig. Für die
Bezeichnung gilt dieselbe Einschränkung wie bei den
Investitionspositionen: Sie muss nichtleer sein und darf keine der
reservierten Spaltenbezeichnungen der Cashflow-Zeitreihe (etwa
\texttt{opex\_gesamt\_eur} oder \texttt{erloes\_eur}) tragen, da jede
Position als eigene Spalte in die Zeitreihe geschrieben wird.

\hypertarget{produktionsbasierte-positionen}{%
\subsection{8.3 Produktionsbasierte
Positionen}\label{produktionsbasierte-positionen}}

\textbf{Gemeindeabgabe} (je erzeugter kWh an die Standortgemeinde):

\[ C^{\mathrm{gem}}_t = E_t \cdot c_{\mathrm{gem}} \cdot \Theta_t \]

\textbf{Direktvermarktungskosten} (Bilanzkreis, Prognose, Marktzugang)
in drei Modi:

Modus \textbf{absolut} -- fester Satz je kWh:

\[ C^{\mathrm{dv}}_t = E_t \cdot c_{\mathrm{dv}} \cdot \Theta_t \]

Modus \textbf{Anteil am Großhandelspreis} -- Anteil \(\varphi\) am
nominalen Baseload-Preis \(b_t\) des Szenarios:

\[ C^{\mathrm{dv}}_t = E_t \cdot \frac{b_t}{100} \cdot \varphi \]

Modus \textbf{Anteil am Marktwert Solar} -- derselbe Anteil, bezogen auf
den technologiespezifischen Marktwert:

\[ C^{\mathrm{dv}}_t = E_t \cdot \frac{m_t}{100} \cdot \varphi \]

Marktüblich sind rund 10 \%, bezogen auf den \textbf{Großhandelspreis}:
Der Direktvermarkter rechnet gegen den Spotmarkt ab, nicht gegen den
Marktwert der einzelnen Technologie. Da für Photovoltaik \(m_t < b_t\)
gilt (Kannibalisierung), führen beide relativen Modi bei gleichem
Prozentsatz zu spürbar unterschiedlichen Kosten; der Bezug ist deshalb
eine bewusste Wahl und keine Formalie. Führt das gewählte Szenario keine
Baseload-Kurve, wird ersatzweise \(m_t\) verwendet.

In beiden relativen Modi verändern sich die Kosten proportional zum
Preisniveau. Eine zusätzliche Kostenindexierung erfolgt nicht, da die
Preisentwicklung bereits im nominalen Preis enthalten ist und
andernfalls doppelt berücksichtigt würde.

\hypertarget{pacht}{%
\subsection{8.4 Pacht}\label{pacht}}

Die Pacht wird als eigenständige Kostenposition geführt, da sie abhängig
von Vertragsmodell unterschiedlich bemessen wird.

\textbf{Modus fix} -- fester Betrag je installierter kWp und Jahr:

\[ C^{\mathrm{pacht}}_t = p_{\mathrm{kWp}} \cdot P \cdot \Theta_t \]

\textbf{Modus Umsatzbeteiligung} -- Anteil \(\beta\) am Jahresumsatz,
mindestens aber eine indexierte Mindestpacht je Hektar:

\[ C^{\mathrm{pacht}}_t = \max\left(\beta \cdot R_t,\ \ p_{\mathrm{ha}} \cdot A_{\mathrm{ha}} \cdot \Theta_t\right) \]

Marktüblich sind \(\beta \approx 5{,}5\,\%\). Die Maximumbildung ist
wirtschaftlich relevant: In frühen Jahren dominiert typischerweise die
Umsatzbeteiligung, in späten Jahren die stetig steigende Mindestpacht --
insbesondere nach Auslaufen der Förderung und mit fortschreitender
Degradation. Ist keine Projektfläche gesetzt, wirkt die Mindestpacht wie
null, und es bleibt die reine Umsatzbeteiligung.

\begin{quote}
Die Umsatzbeteiligung verwendet \(R_t\) aus Schritt 3. Diese
Abhängigkeit begründet die Anordnung des Erlösmoduls vor dem
Betriebskostenmodul.
\end{quote}

\hypertarget{summe}{%
\subsection{8.5 Summe}\label{summe}}

\[ C_t = \sum_j C^{(j)}_t + C^{\mathrm{gem}}_t + C^{\mathrm{dv}}_t + C^{\mathrm{pacht}}_t \]

\textbf{Ausgang.} \(C_t\) sowie jede Einzelposition als eigene Spalte --
die Grundlage der aufgeschlüsselten Kostendarstellung in Oberfläche,
Excel- Export und PDF-Bericht.

\hypertarget{schritt-5-finanzierung}{%
\section{9 Schritt 5 -- Finanzierung}\label{schritt-5-finanzierung}}

\textbf{Zweck.} Zins, Tilgung und Restschuld über die Kreditlaufzeit.

\textbf{Codestelle.} \nolinkurl{engine/financing.py},
\texttt{calculate\_financing()}.

\hypertarget{kapitalstruktur}{%
\subsection{9.1 Kapitalstruktur}\label{kapitalstruktur}}

\[ D = I \cdot (1 - e), \qquad EK = I \cdot e \]

Die Kreditsumme ergibt sich residual aus Investitionsvolumen und
Eigenkapitalquote. Das Modell enthält keine gesonderte
Zwischenfinanzierung, keine Bauzeitzinsen und keine Disagio- oder
Gebührenposition; entsprechende Kosten sind gegebenenfalls im CAPEX zu
erfassen.

\hypertarget{konventionen}{%
\subsection{9.2 Konventionen}\label{konventionen}}

\begin{itemize}
\tightlist
\item
  Zinsen eines Jahres fallen auf den \textbf{Jahresanfangsstand} \(B_t\)
  an.
\item
  Tilgung fließt \textbf{nachschüssig} am Jahresende.
\item
  Der Schuldendienst endet nach \(n\) Tilgungsraten (bei tilgungsfreiem
  Anlaufjahr entsprechend ein Jahr später).
\end{itemize}

\[ B_1 = D, \qquad B_{t+1} = \left(B_t - T_t\right)^{+} \]

\hypertarget{zinsaufwand}{%
\subsection{9.3 Zinsaufwand}\label{zinsaufwand}}

\[ Z_t = B_t \cdot i \cdot \left(f \cdot \mathbf{1}_{[t=1]} + \mathbf{1}_{[t>1]}\right) \quad \text{für } t \leq t^{\mathrm{ende}}, \qquad Z_t = 0 \text{ sonst} \]

Dabei ist \(f\) der Zins-Anteilsfaktor des Anlaufjahres aus Abschnitt
5.3 (act/365 oder 30/360) und

\[ t^{\mathrm{ende}} = n + \mathbf{1}_{[\text{tilgungsfreies Anlaufjahr}]} \]

\hypertarget{tilgung}{%
\subsection{9.4 Tilgung}\label{tilgung}}

\textbf{Annuitätentilgung.} Die konstante Rate folgt der
Standard-Annuität (identisch zu \texttt{PMT} in Tabellenkalkulationen):

\[ \mathrm{Ann} = D \cdot \frac{i}{1 - (1+i)^{-n}} \]

\[ T_t = \mathrm{Ann} - Z_t \]

\textbf{Lineare Tilgung.} Konstante Tilgungsrate, fallender
Schuldendienst:

\[ T_t = \frac{D}{n} \]

\textbf{Tilgungsfreies Anlaufjahr.} Ist es aktiviert, gilt \(T_1 = 0\);
die Tilgung beginnt in Jahr 2. Die \textbf{Anzahl} der Raten bleibt
\(n\); der Schuldendienstzeitraum verlängert sich dadurch um ein Jahr.
Da im ersten Jahr keine Tilgung erfolgt, wird der Zins auch im zweiten
Jahr auf die vollständige Kreditsumme berechnet.

\[ T_t = 0 \quad \text{für } t < t^{\mathrm{tilg}}, \qquad t^{\mathrm{tilg}} = 1 + \mathbf{1}_{[\text{tilgungsfreies Anlaufjahr}]} \]

\textbf{Schuldendienst.}

\[ \mathrm{DS}_t = Z_t + T_t \]

\hypertarget{wechselwirkung-mit-dem-unterjuxe4hrigen-anlaufjahr}{%
\subsection{9.5 Wechselwirkung mit dem unterjährigen
Anlaufjahr}\label{wechselwirkung-mit-dem-unterjuxe4hrigen-anlaufjahr}}

Der Anteilsfaktor \(f\) reduziert \textbf{ausschließlich die Zinslast}
des ersten Jahres. Die Tilgung folgt unverändert der Annuitäts- bzw.
Linearformel. Bei Annuitätentilgung bedeutet das: Weil von der fixen
Rate im ersten Jahr weniger Zins abgeht, wird entsprechend mehr getilgt;
das Darlehen kann dadurch geringfügig vor Ablauf der nominellen Laufzeit
vollständig getilgt sein. Dieser Effekt ergibt sich systematisch aus
einem unterjährigen ersten Zinszeitraum und stellt keine numerische
Approximation dar.

Die Kappung \(B_{t+1} = (B_t - T_t)^{+}\) verhindert in jedem Fall einen
negativen Darlehensstand.

\textbf{Ausgang.} \(Z_t\), \(T_t\), \(\mathrm{DS}_t\), \(B_t\)
(Jahresanfang) und \(B_{t+1}\) (Jahresende).

\hypertarget{schritt-6-ertragsteuern}{%
\section{10 Schritt 6 -- Ertragsteuern}\label{schritt-6-ertragsteuern}}

\textbf{Zweck.} Steuerliche Bemessungsgrundlage und Steuerzahlung je
Jahr.

\textbf{Codestelle.} \nolinkurl{engine/tax.py}, \texttt{calculate\_tax()}.

Dieses Modul weist als einziges Teilmodell eine periodenübergreifende
Zustandsfortschreibung auf: Der Verlustvortragsbestand einer Periode
hängt vom Bestand der Vorperiode ab. Die Berechnung erfolgt daher
sequenziell über die Zeitachse.

\hypertarget{ergebnis-vor-abschreibung}{%
\subsection{10.1 Ergebnis vor
Abschreibung}\label{ergebnis-vor-abschreibung}}

\[ \mathrm{EBT}^{\mathrm{vA}}_t = R_t - C_t - Z_t \]

Die Bemessungsgrundlage setzt \textbf{nach} Zinsen an (Zinsen sind
Betriebsausgaben), aber \textbf{vor} Abschreibung -- diese wird im
nächsten Schritt modusabhängig abgezogen. Die Tilgung ist erfolgsneutral
und geht korrekterweise nicht ein.

\hypertarget{abschreibung}{%
\subsection{10.2 Abschreibung}\label{abschreibung}}

\[ A_t = \frac{I}{n_{\mathrm{AfA}}} \cdot \mathbf{1}_{[\,t\, \leq\, n_{\mathrm{AfA}}\,]} \]

Lineare Abschreibung des gesamten Investitionsvolumens über die
steuerliche Nutzungsdauer. Nach Ablauf der Nutzungsdauer ist das
Wirtschaftsgut voll abgeschrieben; eine weitere Abschreibung wäre
unzulässig und ist ausgeschlossen. Im Pauschalmodus gilt \(A_t = 0\).

\hypertarget{modusabhuxe4ngige-parameter}{%
\subsection{10.3 Modusabhängige
Parameter}\label{modusabhuxe4ngige-parameter}}

\begin{longtable}[]{@{}
  >{\raggedright\arraybackslash}p{(\columnwidth - 8\tabcolsep) * \real{0.2000}}
  >{\raggedright\arraybackslash}p{(\columnwidth - 8\tabcolsep) * \real{0.2000}}
  >{\raggedright\arraybackslash}p{(\columnwidth - 8\tabcolsep) * \real{0.2000}}
  >{\raggedright\arraybackslash}p{(\columnwidth - 8\tabcolsep) * \real{0.2000}}
  >{\raggedright\arraybackslash}p{(\columnwidth - 8\tabcolsep) * \real{0.2000}}@{}}
\toprule\noalign{}
\begin{minipage}[b]{\linewidth}\raggedright
Modus
\end{minipage} & \begin{minipage}[b]{\linewidth}\raggedright
\(A_t\)
\end{minipage} & \begin{minipage}[b]{\linewidth}\raggedright
Freibetrag \(\Phi\)
\end{minipage} & \begin{minipage}[b]{\linewidth}\raggedright
Satz \(\tau\)
\end{minipage} & \begin{minipage}[b]{\linewidth}\raggedright
Grenze \(\gamma\)
\end{minipage} \\
\midrule\noalign{}
\endhead
\bottomrule\noalign{}
\endlastfoot
Pauschal auf EBT & 0 & 0 & Eingabewert & Eingabewert \\
Körperschaftsteuer (AT) & linear über \(n_{\mathrm{AfA}}\) & Eingabewert
& Eingabewert & Eingabewert (gesetzlich 0,75) \\
Gewerbesteuer (DE) & linear über \(n_{\mathrm{AfA}}\) & 24.500 € &
\(0{,}035 \cdot H/100\) & 0 \\
\end{longtable}

Der effektive Gewerbesteuersatz folgt der gesetzlichen Systematik aus
Steuermesszahl und gemeindlichem Hebesatz \(H\):

\[ \tau_{\mathrm{GewSt}} = 0{,}035 \cdot \frac{H}{100} \]

Bei einem Hebesatz von 400 \% ergibt das \(\tau = 14{,}0\,\%\).

\hypertarget{steuerliches-ergebnis-vor-verlustvortrag}{%
\subsection{10.4 Steuerliches Ergebnis vor
Verlustvortrag}\label{steuerliches-ergebnis-vor-verlustvortrag}}

\[ G_t = \mathrm{EBT}^{\mathrm{vA}}_t - A_t - \Phi \]

\hypertarget{verlustvortrag}{%
\subsection{10.5 Verlustvortrag}\label{verlustvortrag}}

Nach österreichischem Recht (§ 8 Abs. 4 Z 2 KStG) sind Verluste zeitlich
\textbf{unbegrenzt} vortragbar, in einem Gewinnjahr aber nur bis zur
Verrechnungsgrenze \(\gamma\) (gesetzlich 75 \%) des steuerlichen
Ergebnisses verrechenbar. Der Rest ist in jedem Fall zu versteuern.

Mit dem Vortragsbestand \(V_t\) zu Jahresbeginn (\(V_1 = 0\)):

\[ U_t = \min\left(V_t,\ \gamma \cdot G_t\right) \cdot \mathbf{1}_{[\,G_t\, >\, 0\,]} \]

\[ G^{\mathrm{st}}_t = \left(G_t - U_t\right)^{+} \]

\[ S_t = G^{\mathrm{st}}_t \cdot \tau \]

\[ V_{t+1} = V_t - U_t + \left(-G_t\right)^{+} \]

Dabei ist \(U_t\) der genutzte Vortrag, \(G^{\mathrm{st}}_t\) das
tatsächlich versteuerte Ergebnis und \((-G_t)^{+}\) der im Jahr \(t\)
neu entstandene Verlust.

Ein optionaler Deaktivierungsschalter für den Verlustvortrag ist nicht
vorgesehen, da dessen Berücksichtigung im österreichischen Steuermodus
Bestandteil der implementierten Rechtslogik ist. Steuerung erfolgt
ausschließlich über die Verrechnungsgrenze \(\gamma\); mit
\(\gamma = 0\) ist der Vortrag faktisch deaktiviert (so wird die
deutsche Gewerbesteuer abgebildet).

\begin{quote}
\textbf{Annahme (Gewerbesteuer).} Der gewerbesteuerliche Verlustvortrag
nach § 10a GewStG wird nicht abgebildet: Das als Referenz validierte
Modell (Abgleich mit einer realen Projekt-Excel) berücksichtigt ihn
ebenfalls nicht. Jedes Jahr wird unabhängig betrachtet.
\end{quote}

\hypertarget{ausweis-der-steuerlichen-zwischengruxf6uxdfen}{%
\subsection{10.6 Ausweis der steuerlichen
Zwischengrößen}\label{ausweis-der-steuerlichen-zwischengruxf6uxdfen}}

Neben \(S_t\) werden \(A_t\), \(V_t\), \(G_t\), \(U_t\), \(V_{t+1}\) und
\(G^{\mathrm{st}}_t\) ausgegeben. Der vollständige Ausweis dieser
Zwischengrößen ermöglicht die periodenbezogene Prüfung der Überleitung
vom Ergebnis vor Steuern zur tatsächlichen Steuerzahlung.

\textbf{Ausgang.} \(A_t\), \(V_t\), \(G_t\), \(U_t\), \(V_{t+1}\),
\(G^{\mathrm{st}}_t\), \(S_t\).

\hypertarget{schritt-7-cashflow-rechnung}{%
\section{11 Schritt 7 --
Cashflow-Rechnung}\label{schritt-7-cashflow-rechnung}}

\textbf{Zweck.} Zusammenführung aller Zeitreihen zum Equity-Cashflow.

\textbf{Codestelle.} \nolinkurl{engine/cashflow.py},
\texttt{calculate\_cashflow()}.

Dieses Modul führt die zuvor berechneten Zeitreihen ohne zusätzliche
fachliche Annahmen zusammen. Die Trennung von Berechnung und Aggregation
ermöglicht eine eindeutige Zuordnung etwaiger Abweichungen zum jeweils
verantwortlichen Teilmodell.

\hypertarget{investitionszeitpunkt-t-0}{%
\subsection{\texorpdfstring{11.1 Investitionszeitpunkt
\(t = 0\)}{11.1 Investitionszeitpunkt t = 0}}\label{investitionszeitpunkt-t-0}}

\[ \mathrm{CF}^{\mathrm{op}}_0 = 0, \qquad \mathrm{CF}^{\mathrm{inv}}_0 = -I, \qquad \mathrm{CF}^{\mathrm{fin}}_0 = D \]

\[ \mathrm{CF}_0 = -I + D = -I \cdot e = -EK \]

Der Nettozahlungsmittelabfluss zum Zeitpunkt \(t=0\) entspricht damit
dem \textbf{Eigenkapitaleinsatz}. Die Zeile trägt als Datum das
Inbetriebnahmedatum -- sie ist der Bezugspunkt aller Diskontierungen.

\hypertarget{betriebsjahre-t-geq-1}{%
\subsection{\texorpdfstring{11.2 Betriebsjahre
\(t \geq 1\)}{11.2 Betriebsjahre t \textbackslash geq 1}}\label{betriebsjahre-t-geq-1}}

\[ \mathrm{CF}^{\mathrm{op}}_t = R_t - C_t - Z_t - S_t \]

\[ \mathrm{CF}^{\mathrm{inv}}_t = 0, \qquad \mathrm{CF}^{\mathrm{fin}}_t = -T_t \]

\[ \mathrm{CF}_t = \mathrm{CF}^{\mathrm{op}}_t + \mathrm{CF}^{\mathrm{inv}}_t + \mathrm{CF}^{\mathrm{fin}}_t \]

\[ \mathrm{CF}^{\mathrm{kum}}_t = \sum_{k=0}^{t} \mathrm{CF}_k \]

Der operative Cashflow ist bereits nach Zinsen und Steuern definiert;
die Tilgung erscheint separat im Finanzierungs-Cashflow.
\(\mathrm{CF}_t\) ist damit durchgängig als \textbf{Cashflow aus Sicht
der Eigenkapitalgeber} definiert. Er bildet die gemeinsame
Datengrundlage für XNPV, XIRR und Amortisationszeit.

\hypertarget{schuldendienstdeckungsgrad-dscr}{%
\subsection{11.3 Schuldendienstdeckungsgrad
(DSCR)}\label{schuldendienstdeckungsgrad-dscr}}

\[ \mathrm{CFADS}_t = R_t - C_t - S_t \]

\[ \mathrm{DSCR}_t = \frac{\mathrm{CFADS}_t}{\mathrm{DS}_t} \quad \text{für } \mathrm{DS}_t > 0, \qquad \text{sonst undefiniert} \]

CFADS (\emph{Cash Flow Available for Debt Service}) ist der Cashflow
\textbf{vor} Zinsen. Da die Zinsen bereits als Bestandteil des
Schuldendienstes im Nenner enthalten sind, werden sie im Zähler nicht
erneut abgezogen. Für Perioden ohne Schuldendienst ist der DSCR nicht
definiert. Diese Perioden werden bei der Ermittlung des Minimums und bei
Verlaufsanalysen nicht berücksichtigt.

\hypertarget{ergebnisstruktur}{%
\subsection{11.4 Ergebnisstruktur}\label{ergebnisstruktur}}

Die Cashflow-Zeitreihe führt neben den Aggregaten alle Erklärgrößen mit:
Produktion, Marktwert real und nominal, Vergütungssatz, Erlösaufteilung,
Betriebskosten je Einzelposition, Zinsen, Tilgung, sämtliche
Steuergrößen. Das Spaltenschema wird beim Erzeugen erzwungen -- fehlt
eine Spalte, bricht die Rechnung ab, statt eine unvollständige Tabelle
weiterzureichen.

\textbf{Ausgang.} Vollständige Cashflow-Tabelle, Zeilen
\(t = 0 \dots N\).

\hypertarget{dscr-kovenanten-ausschuxfcttungssperre-und-nachschusspflicht}{%
\subsection{11.5 DSCR-Kovenanten: Ausschüttungssperre und
Nachschusspflicht}\label{dscr-kovenanten-ausschuxfcttungssperre-und-nachschusspflicht}}

\textbf{Codestelle.} \nolinkurl{engine/covenants.py},
\texttt{analysiere\_kovenanten()}.

Kreditverträge belegen den Schuldendienstdeckungsgrad mit zwei
Schwellen. Beide wirken \textbf{nicht} auf die Cashflow-Rechnung zurück;
sie werden auf der fertigen Zeitreihe ausgewertet.

\[ s_{\mathrm{trap}} = \text{Cash-Trap-Schwelle (Vorbelegung } 1{,}10\text{)}, \qquad s_{\mathrm{eod}} = \text{Event-of-Default-Schwelle (Vorbelegung } 1{,}05\text{)} \]

\textbf{Ereignisse.} Für alle Perioden mit \(\mathrm{DS}_t > 0\):

\[ \text{Cash Trap}_t \iff \mathrm{DSCR}_t < s_{\mathrm{trap}}, \qquad \text{Event of Default}_t \iff \mathrm{DSCR}_t < s_{\mathrm{eod}} \]

\textbf{Nachschussbetrag.} Bei einem Event of Default wird der Verstoß
üblicherweise durch eine Eigenkapitaleinlage geheilt (\emph{Equity
Cure}), und zwar in der Höhe, die den Deckungsgrad gerade wieder auf die
Schwelle hebt. Zusätzlich ist eine reine Zahlungslücke stets zu decken
-- auch dann, wenn die Schwelle so niedrig gesetzt ist, dass sie formal
nicht greift:

\[ N_t = \max\left( \mathbf{1}_{[\mathrm{DSCR}_t\, <\, s_{\mathrm{eod}}]} \left(s_{\mathrm{eod}} \mathrm{DS}_t - \mathrm{CFADS}_t\right)^{+},\ \left(-\mathrm{CF}_t\right)^{+} \right) \]

\textbf{Ausschüttung und Reserve.} Der Cash Trap sperrt die
Ausschüttung; der freie Cashflow verbleibt dann als Reserve in der
Gesellschaft:

\[ \text{Ausschüttung } G_t^{\mathrm{aus}} = \left(\mathrm{CF}_t\right)^{+} \cdot \mathbf{1}_{[\mathrm{DSCR}_t\, \geq\, s_{\mathrm{trap}}]}, \qquad \text{Reservezugang } \left(\mathrm{CF}_t\right)^{+} \cdot \mathbf{1}_{[\mathrm{DSCR}_t\, <\, s_{\mathrm{trap}}]} \]

\textbf{Deckungswasserfall.} Der Nachschuss wird in fester Reihenfolge
aus drei Quellen gedeckt. Mit dem Reservebestand \(Q_t\) und dem Bestand
bereits ausgeschütteter, rückführbarer Mittel \(A_t\):

\[ n^{\mathrm{res}}_t = \min(N_t,\ Q_t), \qquad n^{\mathrm{aus}}_t = \min\left(N_t - n^{\mathrm{res}}_t,\ A_t\right), \qquad n^{\mathrm{ext}}_t = N_t - n^{\mathrm{res}}_t - n^{\mathrm{aus}}_t \]

\[ Q_{t+1} = Q_t - n^{\mathrm{res}}_t + \left(\mathrm{CF}_t\right)^{+}\mathbf{1}_{[\mathrm{DSCR}_t\, <\, s_{\mathrm{trap}}]}, \qquad A_{t+1} = A_t - n^{\mathrm{aus}}_t + G_t^{\mathrm{aus}} \]

Die ersten beiden Quellen sind Mittel, die das Projekt zuvor
\textbf{selbst erwirtschaftet} hat: einbehaltener Cashflow und bereits
an die Gesellschafter ausgekehrtes Kapital. Nur der Rest erfordert
\textbf{zusätzliches externes Kapital}:

\[ N^{\mathrm{ext}} = \sum_{t=1}^{N} n^{\mathrm{ext}}_t > 0 \quad \Longleftrightarrow \quad \text{externe Kapitalzuführung erforderlich} \]

Diese Unterscheidung ist die eigentliche Aussage der Prüfung: Ein
Nachschussbedarf, der aus eigener Kraft gedeckt werden kann, ist ein
Liquiditäts-, kein Finanzierungsproblem.

\begin{quote}
\textbf{Annahme.} Der Deckungswasserfall unterstellt, dass
ausgeschüttete Mittel bei Bedarf in voller Höhe zurückgeführt werden
können. Er beziffert damit die Obergrenze der aus eigener Kraft
möglichen Deckung; ob die Gesellschafter tatsächlich zurückführen, ist
eine Frage des Gesellschaftsvertrags und nicht Gegenstand des Modells.
\end{quote}

\hypertarget{schritt-8-bewertungskennzahlen}{%
\section{12 Schritt 8 --
Bewertungskennzahlen}\label{schritt-8-bewertungskennzahlen}}

\textbf{Zweck.} Ableitung der entscheidungsrelevanten Barwert-,
Rendite-, Amortisations- und Finanzierungskennzahlen aus der datierten
Cashflow-Zeitreihe.

\textbf{Codestelle.} \nolinkurl{engine/kpis.py} und
\nolinkurl{engine/analytics.py} (\texttt{calculate\_lcoe}).

\hypertarget{nettobarwert-bei-taggenauer-diskontierung-xnpv}{%
\subsection{12.1 Nettobarwert bei taggenauer Diskontierung
(XNPV)}\label{nettobarwert-bei-taggenauer-diskontierung-xnpv}}

Die in Abschnitt 2.2 eingeführte Barwertfunktion wird auf Act/365-Basis
taggenau ausgewertet. Sie entspricht damit der Berechnungslogik der
Tabellenkalkulationsfunktion \texttt{XNPV}:

\[ \operatorname{XNPV}(r) = \sum_{t=0}^{N} \frac{\mathrm{CF}_t}{(1+r)^{\frac{\Delta_t}{365}}}, \qquad \Delta_t = \text{Tage}(d_0,\ d_t) \]

Dabei ist \(d_0\) das Investitionsdatum (Inbetriebnahmedatum) und
\(d_t\) das Ende des Betriebsjahres \(t\). Die jahresbasierte
Vereinfachung \((1+r)^{-t}\) wird nicht verwendet, da sie bei
unterjähriger Inbetriebnahme systematisch von der taggenauen
Diskontierung abweicht.

Der ausgewiesene Nettobarwert entspricht \(\operatorname{XNPV}(r)\) bei
einem frei wählbaren Diskontsatz (Vorbelegung 8 \%). Die zugehörige
Barwertkurve wertet dieselbe Funktion an 21 Stützstellen zwischen 0 \%
und 10 \% in Schritten von 0,5 Prozentpunkten aus. Ein Nulldurchgang
dieser Funktion entspricht einem internen Zinsfuß.

\hypertarget{interner-zinsfuuxdf-bei-unregelmuxe4uxdfigen-zahlungszeitpunkten-xirr}{%
\subsection{12.2 Interner Zinsfuß bei unregelmäßigen Zahlungszeitpunkten
(XIRR)}\label{interner-zinsfuuxdf-bei-unregelmuxe4uxdfigen-zahlungszeitpunkten-xirr}}

Der Equity-IRR ist gemäß Abschnitt 2.2.2 als Nullstelle der
XNPV-Funktion definiert:

\[ \operatorname{XNPV}(r^{*}) = 0 \]

Gelöst wird mit dem Verfahren von Brent (\texttt{brentq}), einer
Kombination aus Bisektion, Sekantenverfahren und inverser quadratischer
Interpolation. Das Verfahren benötigt ein Intervall mit
Vorzeichenwechsel. Um auch extreme Projekte abzudecken, wird das
Suchintervall schrittweise erweitert:

\[ \left[-0{,}9999,\ 10\right] \ \rightarrow\ \left[-0{,}9999,\ 100\right] \ \rightarrow\ \left[-0{,}9999,\ 1000\right] \]

Vorab wird geprüft, ob der Cashflow überhaupt einen Vorzeichenwechsel
enthält. Fehlt ein Vorzeichenwechsel, ist der interne Zinsfuß für die
vorliegende Cashflow-Folge nicht bestimmbar. In diesem Fall wird kein
numerischer Ersatzwert ausgewiesen.

\begin{quote}
Die Mehrdeutigkeit des IRR bei mehrfachem Vorzeichenwechsel ist eine
bekannte Eigenschaft der Kennzahl. Bei der hier vorliegenden
Cashflow-Struktur (ein Abfluss zu Beginn, danach überwiegend Zuflüsse)
tritt sie praktisch nicht auf; ausgewiesen wird die vom Verfahren
gefundene Nullstelle im Suchintervall.
\end{quote}

\hypertarget{amortisationszeit}{%
\subsection{12.3 Amortisationszeit}\label{amortisationszeit}}

\[ t^{\mathrm{PB}} = \min\left\{\,t \ :\ \mathrm{CF}^{\mathrm{kum}}_t \geq 0\,\right\} \]

Die Amortisation ist \textbf{undiskontiert} definiert (einfache Payback-
Periode) und wird in vollen Betriebsjahren ausgewiesen. Erreicht der
kumulierte Cashflow innerhalb der Betrachtungsdauer nie null, wird kein
Wert ausgewiesen.

\hypertarget{weitere-kennzahlen}{%
\subsection{12.4 Weitere Kennzahlen}\label{weitere-kennzahlen}}

\[ EK = -\mathrm{CF}_0, \qquad I = -\mathrm{CF}^{\mathrm{inv}}_0 \]

\[ \mathrm{DSCR}^{\min} = \min_{t\,:\,\mathrm{DS}_t > 0} \mathrm{DSCR}_t \]

\[ \text{Spezifisches Invest} = \frac{I}{P} \quad \left[\text{€/kWp}\right] \]

\textbf{Equity Value.} Der Nettobarwert enthält den Eigenkapitaleinsatz
des Jahres 0 als Abfluss. Rechnet man ihn wieder hinzu, verbleibt der
Barwert der künftigen Eigenkapital-Cashflows -- der Wert des
Eigenkapitals zum Bewertungsstichtag:

\[ V^{\mathrm{EK}}(r) = \operatorname{XNPV}(r) + EK = \sum_{t=1}^{N} \frac{\mathrm{CF}_t}{(1+r)^{\Delta_t/365}} \]

\textbf{Enterprise Value.} Zuzüglich des zum Stichtag aufgenommenen
Fremdkapitals ergibt sich der Gesamtunternehmenswert:

\[ V^{\mathrm{GK}}(r) = V^{\mathrm{EK}}(r) + D, \qquad D = I - EK \]

\hypertarget{stromgestehungskosten-lcoe}{%
\subsection{12.5 Stromgestehungskosten
(LCOE)}\label{stromgestehungskosten-lcoe}}

\[ v_t = (1 + r)^{-\frac{\Delta_t}{365}} \]

\[ \mathrm{LCOE} = \frac{\sum_{t=0}^{N} \left(-\mathrm{CF}^{\mathrm{inv}}_t + C_t\right) \cdot v_t}{\sum_{t=0}^{N} E_t \cdot v_t} \cdot 100 \quad \left[\text{ct/kWh}\right] \]

Zähler sind die diskontierten Vollkosten (Investition plus
Betriebskosten), Nenner die diskontierte Strommenge; der Faktor 100
wandelt €/kWh in ct/kWh. Die Diskontierung ist dieselbe taggenaue
Act/365-Systematik wie beim XNPV.

Bewusst \textbf{nicht} enthalten sind Zinsen und Steuern: Die
Finanzierungskosten stecken definitionsgemäß im Diskontsatz \(r\); sie
zusätzlich im Zähler zu führen, würde sie doppelt zählen. Der LCOE ist
damit die Kostenseite des Projekts unabhängig von der konkreten
Finanzierungsstruktur und direkt mit Marktpreisen und Zuschlagswerten
vergleichbar.

\textbf{Ausgang.} Equity-IRR, NPV, Payback, Investitionsvolumen,
Eigenkapitaleinsatz, minimaler DSCR, LCOE.

\hypertarget{numerisches-anwendungsbeispiel}{%
\section{13 Numerisches
Anwendungsbeispiel}\label{numerisches-anwendungsbeispiel}}

Dieses Kapitel dokumentiert die numerische Reproduktion eines
vollständig spezifizierten Beispielprojekts. Grundlage ist das
mitgelieferte Beispielprojekt \nolinkurl{data/projects/template-agri.yaml}
mit den globalen Annahmen aus \nolinkurl{data/global_assumptions.yaml}.
Alle Tabellen dieses Kapitels erzeugt
\texttt{python }\nolinkurl{docs/rechenmodell/beispiel.py};
\nolinkurl{tests/test_dokumentation.py} prüft, dass die abgedruckten
Zahlen weiterhin dem Rechenergebnis entsprechen.

\hypertarget{eingangsgruxf6uxdfen}{%
\subsection{13.1 Eingangsgrößen}\label{eingangsgruxf6uxdfen}}

\begin{longtable}[]{@{}
  >{\raggedright\arraybackslash}p{(\columnwidth - 2\tabcolsep) * \real{0.5000}}
  >{\raggedright\arraybackslash}p{(\columnwidth - 2\tabcolsep) * \real{0.5000}}@{}}
\toprule\noalign{}
\begin{minipage}[b]{\linewidth}\raggedright
Größe
\end{minipage} & \begin{minipage}[b]{\linewidth}\raggedright
Wert
\end{minipage} \\
\midrule\noalign{}
\endhead
\bottomrule\noalign{}
\endlastfoot
Nennleistung \(P\) & 3.800 kWp \\
Spezifischer Ertrag \(h\) & 1.400 kWh/kWp \\
Inbetriebnahme & 01/2027 \\
Investitionsvolumen \(I\) & 2.915.100 € \\
Eigenkapitalquote \(e\) & 20,0 \% \\
Fremdkapitalzins \(i\) & 4,20 \% \\
Kreditlaufzeit \(n\) / Tilgungsart & 20 Jahre / Annuität \\
EAG-Zuschlagswert \(z\) (Agri-PV, ohne Abschlag) & 6,50 ct/kWh \\
Förderdauer \(F\) / Betriebsdauer \(N\) & 20 / 30 Jahre \\
Degradation \(d\) / Sicherheitsabschlag \(\sigma\) & 0,25 \%/a / 0 \% \\
Marktpreisszenario & Aurora 10/25 \\
Marktpreisinflation \(\iota\) / Basisjahr \(y_B\) & 2,0 \% / 2025 \\
Kosteninflation \(\kappa\) & 2,0 \% \\
Steuermodus & Körperschaftsteuer, \(\tau = 23\,\%\),
\(n_{\mathrm{AfA}} = 20\) \\
Negativstunden & 6-Stunden-Regel, Modus „Marktwert``, Gewichtung 100
\% \\
\end{longtable}

Standard-Betriebskosten (global, je kWp und Jahr, je 2 \% indexiert ab
Jahr 1): technische Betriebsführung 3,00 €, Wartung/Rücklage
Wechselrichter 1,00 €, Versicherungen 1,00 €, kaufmännische
Betriebsführung 3,00 €, Sonstiges 1,00 € -- zusammen 9,00 €/kWp/Jahr.
Projektspezifisch: Pacht 5,263158 €/kWp/Jahr (fix), Gemeindeabgabe 2,00
€/MWh, Direktvermarktung 1,00 €/MWh.

\hypertarget{betriebsjahr-1-schritt-fuxfcr-schritt}{%
\subsection{13.2 Betriebsjahr 1 Schritt für
Schritt}\label{betriebsjahr-1-schritt-fuxfcr-schritt}}

\textbf{Schritt 1 -- Zeitachse.} Inbetriebnahme am 1. Januar 2027, also
\(\pi_1 = 1\) und \(f = 1\); Kalenderjahr \(y_1 = 2027\).

\textbf{Schritt 2 -- Menge.}

\[ E_1 = 3800 \cdot 1400 \cdot (1-0{,}0025)^0 \cdot 1 \cdot 1 = 5.320.000\ \mathrm{kWh} \]

\textbf{Schritt 3 -- Erlös.} Kurvenwert des Szenarios für 2027:
\(m^{\mathrm{real}}_1 = 4{,}133\) ct/kWh; Negativmengenanteil
\(\nu_1 = 0{,}222\).

\[ m_1 = 4{,}133 \cdot 1{,}02^{\,2027-2025} = 4{,}133 \cdot 1{,}0404 = 4{,}300\ \text{ct/kWh} \]

\[ p_1 = (6{,}50 - 4{,}300)^{+} = 2{,}200\ \text{ct/kWh}, \qquad s_1 = 4{,}300 + 2{,}200 = 6{,}500\ \text{ct/kWh} \]

Modus „Marktwert``: Der Markterlös bleibt für die gesamte Menge
erhalten, nur die Prämie entfällt für den Anteil \(\nu_1\).

\[ R^{\mathrm{markt}}_1 = \frac{5.320.000 \cdot 4{,}300}{100} = 228.759\ \text{€} \]

\[ R^{\text{Prämie}}_1 = \frac{5.320.000 \cdot (1 - 0{,}222) \cdot 2{,}200}{100} = 91.058\ \text{€} \]

\[ R_1 = 319.817\ \text{€} \]

\textbf{Schritt 4 -- Betriebskosten.} Im ersten Jahr ist
\(\Theta_1 = 1{,}02^0 = 1\), alle Indexexponenten sind null:

\begin{longtable}[]{@{}lll@{}}
\toprule\noalign{}
Position & Rechnung & Betrag \\
\midrule\noalign{}
\endhead
\bottomrule\noalign{}
\endlastfoot
Standardpositionen & \(9{,}00 \cdot 3800\) & 34.200 € \\
Pacht (fix) & \(5{,}263158 \cdot 3800\) & 20.000 € \\
Gemeindeabgabe & \(5.320.000 \cdot 0{,}002\) & 10.640 € \\
Direktvermarktung & \(5.320.000 \cdot 0{,}001\) & 5.320 € \\
\textbf{Summe \(C_1\)} & & \textbf{70.160 €} \\
\end{longtable}

\textbf{Schritt 5 -- Finanzierung.}

\[ D = 2.915.100 \cdot 0{,}8 = 2.332.080\ \text{€} \]

\[ Z_1 = 2.332.080 \cdot 0{,}042 \cdot 1 = 97.947\ \text{€} \]

\[ \mathrm{Ann} = 2.332.080 \cdot \frac{0{,}042}{1 - 1{,}042^{-20}} = \frac{97.947}{0{,}56083} = 174.651\ \text{€} \]

\[ T_1 = 174.651 - 97.947 = 76.704\ \text{€} \]

\textbf{Schritt 6 -- Steuern.}

\[ \mathrm{EBT}^{\mathrm{vA}}_1 = 319.817 - 70.160 - 97.947 = 151.710\ \text{€} \]

\[ A_1 = \frac{2.915.100}{20} = 145.755\ \text{€} \]

\[ G_1 = 151.710 - 145.755 - 0 = 5.955\ \text{€} \]

Kein Vortragsbestand (\(V_1 = 0\)), also \(U_1 = 0\) und
\(G^{\mathrm{st}}_1 = 5.955\) €:

\[ S_1 = 5.955 \cdot 0{,}23 = 1.370\ \text{€} \]

\textbf{Schritt 7 -- Cashflow.}

\[ \mathrm{CF}^{\mathrm{op}}_1 = 319.817 - 70.160 - 97.947 - 1.370 = 150.340\ \text{€} \]

\[ \mathrm{CF}_1 = 150.340 - 76.704 = 73.636\ \text{€} \]

\[ \mathrm{CFADS}_1 = 319.817 - 70.160 - 1.370 = 248.287\ \text{€}, \qquad \mathrm{DSCR}_1 = \frac{248.287}{174.651} = 1{,}42 \]

\textbf{Jahr 0.}

\[ \mathrm{CF}_0 = -2.915.100 + 2.332.080 = -583.020\ \text{€} \]

\hypertarget{ergebniszeitreihe-ausgewuxe4hlte-jahre}{%
\subsection{13.3 Ergebniszeitreihe (ausgewählte
Jahre)}\label{ergebniszeitreihe-ausgewuxe4hlte-jahre}}

\begin{longtable}[]{@{}
  >{\raggedright\arraybackslash}p{(\columnwidth - 18\tabcolsep) * \real{0.1000}}
  >{\raggedright\arraybackslash}p{(\columnwidth - 18\tabcolsep) * \real{0.1000}}
  >{\raggedright\arraybackslash}p{(\columnwidth - 18\tabcolsep) * \real{0.1000}}
  >{\raggedright\arraybackslash}p{(\columnwidth - 18\tabcolsep) * \real{0.1000}}
  >{\raggedright\arraybackslash}p{(\columnwidth - 18\tabcolsep) * \real{0.1000}}
  >{\raggedright\arraybackslash}p{(\columnwidth - 18\tabcolsep) * \real{0.1000}}
  >{\raggedright\arraybackslash}p{(\columnwidth - 18\tabcolsep) * \real{0.1000}}
  >{\raggedright\arraybackslash}p{(\columnwidth - 18\tabcolsep) * \real{0.1000}}
  >{\raggedright\arraybackslash}p{(\columnwidth - 18\tabcolsep) * \real{0.1000}}
  >{\raggedright\arraybackslash}p{(\columnwidth - 18\tabcolsep) * \real{0.1000}}@{}}
\toprule\noalign{}
\begin{minipage}[b]{\linewidth}\raggedright
Jahr
\end{minipage} & \begin{minipage}[b]{\linewidth}\raggedright
Ertrag (kWh)
\end{minipage} & \begin{minipage}[b]{\linewidth}\raggedright
Marktwert (ct/kWh)
\end{minipage} & \begin{minipage}[b]{\linewidth}\raggedright
Vergütung (ct/kWh)
\end{minipage} & \begin{minipage}[b]{\linewidth}\raggedright
Erlös (€)
\end{minipage} & \begin{minipage}[b]{\linewidth}\raggedright
OPEX (€)
\end{minipage} & \begin{minipage}[b]{\linewidth}\raggedright
Zinsen (€)
\end{minipage} & \begin{minipage}[b]{\linewidth}\raggedright
Tilgung (€)
\end{minipage} & \begin{minipage}[b]{\linewidth}\raggedright
Steuer (€)
\end{minipage} & \begin{minipage}[b]{\linewidth}\raggedright
Equity-CF (€)
\end{minipage} \\
\midrule\noalign{}
\endhead
\bottomrule\noalign{}
\endlastfoot
1 & 5.320.000 & 4,300 & 6,500 & 319.817 & 70.160 & 97.947 & 76.704 &
1.370 & 73.636 \\
2 & 5.306.700 & 4,470 & 6,500 & 322.311 & 71.523 & 94.726 & 79.925 &
2.371 & 73.766 \\
3 & 5.293.433 & 4,749 & 6,500 & 329.055 & 72.912 & 91.369 & 83.282 &
4.374 & 77.117 \\
20 & 5.072.906 & 7,304 & 7,304 & 370.525 & 101.130 & 7.040 & 167.612 &
26.818 & 67.926 \\
21 & 5.060.224 & 7,447 & 7,447 & 376.834 & 103.096 & 0 & 0 & 62.960 &
210.778 \\
30 & 4.947.501 & 8,142 & 8,142 & 402.842 & 122.609 & 0 & 0 & 64.454 &
215.779 \\
\end{longtable}

Drei Muster lassen sich daran ablesen:

\begin{itemize}
\tightlist
\item
  \textbf{Bis Jahr 20} liegt der nominale Marktwert unter dem
  Zuschlagswert von 6,50 ct/kWh; der Vergütungssatz ist konstant 6,50,
  die Differenz wird als Marktprämie gezahlt. Ab Jahr 20 hat die
  Inflationierung den Marktwert über den Zuschlagswert gehoben -- die
  Prämie ist von selbst auf null gelaufen, noch bevor die Förderdauer
  endet.
\item
  \textbf{Ab Jahr 21} entfallen Zins und Tilgung (Kreditlaufzeit 20
  Jahre). Der Equity-Cashflow springt von rund 68.000 € auf rund 211.000
  €, gleichzeitig steigt die Steuerlast deutlich, weil die AfA
  (ebenfalls 20 Jahre) ausgelaufen ist.
\item
  \textbf{Die Steuerlast der ersten Jahre ist klein}, weil AfA (145.755
  €/a) und Zinsen den Großteil des Ergebnisses aufzehren.
\end{itemize}

\hypertarget{kennzahlen}{%
\subsection{13.4 Kennzahlen}\label{kennzahlen}}

\begin{longtable}[]{@{}ll@{}}
\toprule\noalign{}
Kennzahl & Wert \\
\midrule\noalign{}
\endhead
\bottomrule\noalign{}
\endlastfoot
Investitionsvolumen \(I\) & 2.915.100 € \\
Eigenkapitaleinsatz \(EK\) (Jahr 0) & 583.020 € \\
EK-Rendite (XIRR) & 13,80 \% \\
NPV bei 8 \% & 442.171 € \\
Minimaler DSCR & 1,37 \\
Payback (kumulierter Equity-CF \(\geq 0\)) & Jahr 8 \\
Summe Erlöse über 30 Jahre & 10.725.258 € \\
\end{longtable}

\hypertarget{sensitivituxe4t-risiko-und-gebotsuntergrenze}{%
\section{14 Sensitivität, Risiko und
Gebotsuntergrenze}\label{sensitivituxe4t-risiko-und-gebotsuntergrenze}}

\textbf{Codestelle.} \nolinkurl{engine/sensitivity.py} und
\nolinkurl{engine/analytics.py}.

Alle Auswertungen dieses Kapitels arbeiten ausschließlich über den
aufgelösten Parametersatz: Sie mutieren \texttt{EffectiveAssumptions}
und rufen die vollständige Bewertungskette erneut auf. Sie kennen die
internen Rechenmodule nicht und bleiben deshalb bei Änderungen an der
Engine automatisch konsistent.

\hypertarget{treibermutationen}{%
\subsection{14.1 Treibermutationen}\label{treibermutationen}}

Ein Treiber \(\theta\) wird multiplikativ um den Faktor \(\lambda\)
variiert:

\begin{longtable}[]{@{}
  >{\raggedright\arraybackslash}p{(\columnwidth - 2\tabcolsep) * \real{0.5000}}
  >{\raggedright\arraybackslash}p{(\columnwidth - 2\tabcolsep) * \real{0.5000}}@{}}
\toprule\noalign{}
\begin{minipage}[b]{\linewidth}\raggedright
Treiber
\end{minipage} & \begin{minipage}[b]{\linewidth}\raggedright
Mutation
\end{minipage} \\
\midrule\noalign{}
\endhead
\bottomrule\noalign{}
\endlastfoot
EAG-Zuschlagswert & \(z \rightarrow z \cdot \lambda\) \\
Marktwert-Niveau &
\(m^{\mathrm{real}}_y \rightarrow m^{\mathrm{real}}_y \cdot \lambda\)
für alle Stützjahre \(y\) \\
Spezifischer Ertrag & \(h \rightarrow h \cdot \lambda\) \\
Investitionskosten & \(I \rightarrow I \cdot \lambda\) \\
Betriebskosten & \(w_j \rightarrow w_j \cdot \lambda\) für alle
Standardpositionen \\
Pacht & \(p \rightarrow p \cdot \lambda\); bei Umsatzbeteiligung werden
Beteiligungssatz und Mindestpacht gemeinsam skaliert \\
Fremdkapitalzins & \(i \rightarrow i \cdot \lambda\) \\
Negativmengen &
\(\nu^{\mathrm{roh}}_y \rightarrow \nu^{\mathrm{roh}}_y \cdot \lambda\) \\
\end{longtable}

Die Pacht wird in \nolinkurl{engine/opex.py} getrennt von den
Standardpositionen geführt (die Berechnung hängt vom Pachtmodus ab); der
Treiber „Betriebskosten'' enthält sie deshalb nicht, die beiden Treiber
überschneiden sich nicht. Bei Umsatzbeteiligung ist die Zahlung das
Maximum aus Beteiligung und Mindestpacht -- beide Terme werden gemeinsam
skaliert, da eine Variation nur eines Terms wirkungslos bliebe, sobald
der andere führt.

Die Parametervariation setzt auf Ebene der \textbf{Eingangskurve} und
nicht auf Ebene des bereits berechneten Ergebnisses an: Eine Skalierung
des Marktwert-Niveaus verändert dadurch auch die Prämienhöhe, den
Zeitpunkt, an dem der Marktwert den Zuschlagswert übersteigt, und im
relativen Modus die Direktvermarktungskosten. Damit wird der
vollständige fachliche Wirkungszusammenhang berücksichtigt.

\hypertarget{eag-zuschlag-sensitivituxe4t}{%
\subsection{14.2
EAG-Zuschlag-Sensitivität}\label{eag-zuschlag-sensitivituxe4t}}

Fünf Varianten des gebotenen Zuschlagswertes:

\[ z_k = z_{\mathrm{Gebot}} \cdot (1 + \Delta_k), \qquad \Delta_k \in \{+5\,\%,\ +2{,}5\,\%,\ 0,\ -2{,}5\,\%,\ -5\,\%\} \]

Für jede Variante wird die vollständige Bewertung neu gerechnet und IRR
und NPV ausgewiesen. Der Konventionell-Abschlag wirkt weiterhin, da die
Variante am gebotenen Wert ansetzt.

Die Stufen sind bewusst eng gefasst: In einer Ausschreibungsrunde bewegt
sich der Zuschlagswert um wenige Zehntel Cent je Kilowattstunde. Eine
Variation um 10 \% beschreibt keine Entscheidung mehr, die tatsächlich
zur Wahl steht; die weiteren Spannen bleiben der Tornado-Analyse
(Abschnitt 14.3) und der Heatmap (Abschnitt 14.4) vorbehalten.

\hypertarget{tornado-analyse}{%
\subsection{14.3 Tornado-Analyse}\label{tornado-analyse}}

Für jeden der acht Treiber wird die EK-Rendite bei
\(\lambda = 1 - \delta\) und \(\lambda = 1 + \delta\) berechnet
(Vorbelegung \(\delta = 10\,\%\)):

\[ \mathrm{Spanne}(\theta) = \left|\,\mathrm{IRR}(\theta \cdot (1+\delta)) - \mathrm{IRR}(\theta \cdot (1-\delta))\,\right| \]

Die Darstellung wird aufsteigend nach der Ergebnisbandbreite sortiert.
Das Tornado-Bild der Projektfinanzierung. Es beantwortet die Frage,
welcher Parameter die Rendite am stärksten bewegt, und damit, wo
Genauigkeit in der Datenbeschaffung am meisten wert ist.

\hypertarget{irr-heatmap}{%
\subsection{14.4 IRR-Heatmap}\label{irr-heatmap}}

Zwei frei wählbare Treiber werden über ein Raster kombiniert. Für jede
Achse ist eine treiberspezifische Spanne hinterlegt (z. B. ±20 \% beim
Zuschlagswert, ±10 \% beim Ertrag, ±30 \% beim Zins):

\[ \lambda^{x}_a = 1 - \Delta_x + \frac{2\Delta_x (a-1)}{S-1}, \qquad a = 1 \dots S \]

Analog für die \(y\)-Achse; \(S\) ist die Stufenzahl (Vorbelegung 7).
Für jede Rasterzelle \((a,b)\) wird die Bewertung vollständig neu
gerechnet:

\[ \mathrm{IRR}_{a,b} = \mathrm{IRR}\left(\text{Mutation}(\lambda^{x}_a) \circ \text{Mutation}(\lambda^{y}_b)\right) \]

Das sind \(S^2\) vollständige Bewertungen (Vorbelegung 49).

\hypertarget{monte-carlo-simulation}{%
\subsection{14.5 Monte-Carlo-Simulation}\label{monte-carlo-simulation}}

Je Lauf \(k = 1 \dots K\) werden für die aktiven Treiber unabhängige
multiplikative Faktoren gezogen:

\[ \lambda^{(k)}_\theta = \max\left(\,\xi,\ 0{,}4\,\right), \qquad \xi \sim \mathcal{N}\left(1,\ \sigma_\theta^2\right) \]

Vorbelegte Standardabweichungen: Ertrag 5 \%, Marktwert-Niveau 10 \%,
CAPEX 5 \%, OPEX 5 \%. Die Untergrenze 0,4 verhindert unphysikalische
Ausreißer (negative Erträge oder negative Investitionskosten) bei großen
Streuungen.

Gesammelt werden je Lauf IRR, NPV zum gewählten Diskontsatz und der
gesamte Pfad des kumulierten Equity-Cashflows. Daraus entstehen:

\[ \mathrm{P}q_t = \text{empirisches } q\text{-Quantil von } \left\{\mathrm{CF}^{\mathrm{kum},(k)}_t\right\}_{k=1}^{K}, \qquad q \in \{10, 25, 50, 75, 90\} \]

\[ \Pr\left[\mathrm{IRR} \geq \mathrm{Ziel}\right] = \frac{\left|\left\{k \ :\ \mathrm{IRR}^{(k)} \geq \mathrm{Ziel}\right\}\right|}{K} \]

Zwei Festlegungen sind wichtig für die Belastbarkeit:

\begin{itemize}
\tightlist
\item
  Der Zufallsgenerator ist mit einem \textbf{festen Startwert}
  initialisiert. Dieselben Eingaben liefern damit exakt dieselbe
  Verteilung -- Voraussetzung für Caching und für die
  Nachvollziehbarkeit in Gremienunterlagen.
\item
  Läufe ohne berechenbaren IRR gehen in den Nenner, aber nicht in den
  Zähler der Erfolgswahrscheinlichkeit ein und werden damit konservativ
  als \textbf{Zielverfehlung} klassifiziert.
\end{itemize}

Optional kann der EAG-Zuschlagswert je Lauf aus der Gebotsverteilung des
Auktionsmodells gezogen werden (Kapitel 15.8). Dann enthält die
IRR-Verteilung zusätzlich die Auktionsunsicherheit:

\[ z^{(k)} = b^{(k)} \cdot \frac{z}{z_{\mathrm{Gebot}}} \]

Der zweite Faktor erhält den Konventionell-Abschlag des Projekts.

\hypertarget{break-even-zuschlagswert-gebotsassistent}{%
\subsection{14.6 Break-even-Zuschlagswert
(Gebotsassistent)}\label{break-even-zuschlagswert-gebotsassistent}}

Gesucht ist der kleinste anzulegende Wert, mit dem eine Ziel-EK-Rendite
gerade noch erreicht wird -- die \textbf{wirtschaftliche Untergrenze}
für ein Auktionsgebot:

\[ \text{Finde } b^{*} \text{ mit } \mathrm{IRR}\left(b^{*}\right) = \mathrm{IRR}^{\mathrm{Ziel}}, \qquad b^{*} \in [0{,}5,\ 15]\ \text{ct/kWh} \]

Gelöst wird wieder mit \texttt{brentq} auf der monoton steigenden
Funktion \(b \mapsto \mathrm{IRR}(b) - \mathrm{IRR}^{\mathrm{Ziel}}\).
Zwei Randfälle werden ausdrücklich behandelt:

\begin{itemize}
\tightlist
\item
  Ist das Ziel bereits am unteren Rand erreicht, trägt der Marktwert das
  Projekt praktisch allein; ausgewiesen wird die untere Suchgrenze.
\item
  Ist das Ziel am oberen Rand nicht erreichbar, wird \textbf{kein} Wert
  ausgewiesen (das Ziel ist mit diesem Projekt nicht erreichbar).
\end{itemize}

Nicht berechenbare IRR-Werte werden in der Nullstellensuche konservativ
als −100 \% behandelt, damit die Suche nicht an einem undefinierten Wert
abbricht.

\hypertarget{szenarienvergleich}{%
\subsection{14.7 Szenarienvergleich}\label{szenarienvergleich}}

Dasselbe Projekt wird über alle hinterlegten Marktpreisszenarien
gerechnet. Ausgetauscht werden je Szenario ausschließlich die
Marktwertkurve und die Negativmengenkurve; alle übrigen Parameter
bleiben identisch. Verglichen werden IRR, NPV und Gesamterlös sowie die
kumulierten Equity-Cashflow-Pfade.

\hypertarget{empirisches-modell-der-eag-ausschreibungen}{%
\section{15 Empirisches Modell der
EAG-Ausschreibungen}\label{empirisches-modell-der-eag-ausschreibungen}}

\textbf{Zweck.} Aus den veröffentlichten Aggregaten vergangener
Ausschreibungsrunden eine Verteilung der Zuschlagswerte schätzen, daraus
die Zuschlagswahrscheinlichkeit eines Gebots ableiten und ein Gebot zur
gewünschten Risikoneigung empfehlen.

\textbf{Codestelle.} \nolinkurl{engine/auktion.py}, Daten in
\nolinkurl{data/ausschreibungen.yaml}.

\hypertarget{fachlicher-rahmen}{%
\subsection{15.1 Fachlicher Rahmen}\label{fachlicher-rahmen}}

Die österreichischen EAG-Marktprämienausschreibungen für Photovoltaik
funktionieren nach \textbf{Pay-as-Bid mit Gebotspreisreihung}: Gebote
werden aufsteigend gereiht und bis zur ausgeschriebenen Menge
bezuschlagt; jeder Gewinner erhält seinen \textbf{eigenen} Gebotswert
als anzulegenden Wert. Eine Preisobergrenze wird per Verordnung
festgelegt; Gebote darüber sind ungültig.

Das Modell nimmt konsequent die \textbf{Price-Taker-Sicht} ein: Ein
einzelner Bieter beeinflusst den Grenzzuschlagswert nicht messbar. Aus
dieser Perspektive ist die Zuschlagsentscheidung von einer exogenen
Zufallsvariable abhängig, dem Grenzzuschlagswert \(p_m\) der nächsten
Runde.

Die Auktionsliteratur beschreibt für Pay-as-Bid-Verfahren ein Muster,
das auch in den vorliegenden Aggregatdaten erkennbar ist: Bieter setzen
ihr Gebot knapp unter den erwarteten Grenzzuschlag („bid shading``). Bei
schwachem Wettbewerb liegen die Gebote nahe der Preisobergrenze; mit
steigendem Wettbewerb sinken sie und verdichten sich.

\hypertarget{datenlage-und-ihre-konsequenzen}{%
\subsection{15.2 Datenlage und ihre
Konsequenzen}\label{datenlage-und-ihre-konsequenzen}}

Veröffentlicht werden je Runde nur \textbf{Aggregate} der bezuschlagten
Gebote: Minimum, mengengewichteter Mittelwert, Maximum sowie
ausgeschriebene und bezuschlagte Menge und die Preisobergrenze.
Einzelgebote werden nicht veröffentlicht.

Daraus folgt unmittelbar die Methodenwahl: Verfahren auf Rohgeboten
(Kerndichteschätzung, Mischverteilungen) scheiden aus. Es bleibt die
Anpassung \textbf{parametrischer, auf \([0, c]\) beschränkter
Verteilungs- familien} über Momenten- und Quantilbedingungen, wobei
\(c\) die Preisobergrenze ist.

\textbf{Regimeerkennung.} Entscheidend ist, ob eine Runde unterzeichnet
oder überzeichnet war:

\[ \text{Runde } j \text{ gilt als unterzeichnet} \quad \Leftrightarrow \quad \max_j \geq c_j - 0{,}02\ \text{ct/kWh} \]

\begin{itemize}
\tightlist
\item
  \textbf{Unterzeichnet}: Der Höchstzuschlag liegt (bis auf Rundung) an
  der Obergrenze. Das Gebotsvolumen hat nicht ausgereicht, um die Grenze
  zu drücken; alle gültigen Gebote wurden bezuschlagt. Die
  veröffentlichten Aggregate beschreiben damit die \textbf{volle}
  Gebotsverteilung, und die Wettbewerbsquote ist direkt beobachtbar:
\end{itemize}

\[ r_j = \frac{\text{bezuschlagte Menge}_j}{\text{ausgeschriebene Menge}_j} \]

\begin{itemize}
\tightlist
\item
  \textbf{Überzeichnet}: Der Höchstzuschlag liegt unter der Obergrenze;
  die ausgeschriebene Menge ist nahezu vollständig ausgeschöpft. Die
  Aggregate beschreiben nur den \textbf{unteren, bezuschlagten Teil} der
  Gebotsverteilung. Das eingereichte Gebotsvolumen wird nicht
  veröffentlicht -- die Wettbewerbsquote ist \textbf{latent} und muss
  geschätzt werden (Abschnitt 15.5).
\end{itemize}

\hypertarget{verteilungsfamilien-auf-0-c}{%
\subsection{\texorpdfstring{15.3 Verteilungsfamilien auf
\([0, c]\)}{15.3 Verteilungsfamilien auf {[}0, c{]}}}\label{verteilungsfamilien-auf-0-c}}

Alle Familien werden einheitlich über zwei Parameter beschrieben, damit
sie unmittelbar vergleichbar sind und der Wettbewerbs-Link
familienunabhängig formuliert werden kann:

\[ \mu_{\mathrm{rel}} = \frac{\mathbb{E}[b]}{c} \in (0,1) \quad \text{(Lage)}, \qquad \kappa > 0 \quad \text{(Konzentration)} \]

\textbf{Beta (Vorbelegung).} Skalierte Beta-Verteilung auf \([0,c]\):

\[ \frac{b}{c} \sim \mathrm{Beta}(\alpha, \beta), \qquad \alpha = \mu_{\mathrm{rel}} \kappa, \quad \beta = (1-\mu_{\mathrm{rel}})\kappa \]

\[ \mathbb{E}\left[\frac{b}{c}\right] = \frac{\alpha}{\alpha+\beta} = \mu_{\mathrm{rel}}, \qquad \mathrm{Var}\left[\frac{b}{c}\right] = \frac{\mu_{\mathrm{rel}}(1-\mu_{\mathrm{rel}})}{\kappa+1} \]

Die Verteilung ist aufgrund ihres Trägers auf \([0,c]\) beschränkt und
kann unterschiedliche Schiefegrade abbilden. Für \(\mu_{\mathrm{rel}}\)
nahe 1 bei moderatem \(\kappa\) ergibt sich eine Verteilungsform mit
hoher Wahrscheinlichkeitsmasse nahe der Obergrenze und einem
ausgeprägten linken Ausläufer.

\textbf{Kumaraswamy.} Der Beta sehr ähnlich, aber mit analytischer
Quantilfunktion:

\[ F(x) = 1 - \left(1 - \left(\frac{x}{c}\right)^{a}\right)^{\beta}, \qquad Q(q) = c\left(1 - (1-q)^{1/\beta}\right)^{1/a} \]

Die Formparameter heißen hier \(a\) und \(\beta\) (nicht \(a, b\) wie in
der Literatur -- \(b\) ist in diesem Kapitel das Gebot). Ihre Umrechnung
aus \((\mu_{\mathrm{rel}}, \kappa)\) erfolgt numerisch:
\(a = \max(\kappa \mu_{\mathrm{rel}},\ 0{,}05)\), und \(\beta\) wird so
bestimmt, dass die Momentbedingung erfüllt ist:

\[ \mathbb{E}\left[\frac{b}{c}\right] = \beta \cdot B\left(1 + \frac{1}{a},\ \beta\right) = \mu_{\mathrm{rel}} \]

Gelöst wird über \(\ln \beta\) mit \texttt{brentq} im Intervall
\([-8, 10]\).

\textbf{Trunkierte Normalverteilung.} Bewusst als symmetrische
Vergleichsbasis:

\[ b \sim \mathcal{N}\left(\mu_{\mathrm{rel}} c,\ \left(\frac{c}{\kappa}\right)^2\right) \text{ trunkiert auf } [0, c] \]

Sie kann die harte Obergrenze abbilden, aber keinen langen linken
Ausläufer bei gleichzeitig hoher Konzentration rechts ; diese
Einschränkung zeigt sich in der Validierung.

\textbf{Gespiegelte inverse Gamma.} Modelliert den Abstand zur
Obergrenze:

\[ b = c - Y, \qquad Y \sim \mathrm{InvGamma}\left(a,\ \text{scale}\right), \quad a = \max(\kappa,\ 1{,}1) \]

\[ \text{scale} = c\left(1 - \mu_{\mathrm{rel}}\right)(a-1) \ \Longrightarrow\ \mathbb{E}[Y] = c(1-\mu_{\mathrm{rel}}) \ \Longrightarrow\ \mathbb{E}[b] = \mu_{\mathrm{rel}} c \]

Die Dichte fällt zur Obergrenze hin stark gegen null und weist nach
links einen langsam abklingenden Ausläufer auf. Damit lässt sich die für
wettbewerbliche Pay-as-Bid-Verfahren plausible Konzentration der Dichte
knapp unterhalb des erwarteten Grenzzuschlags abbilden. Ihre
prinzipielle Grenze: Masse \textbf{direkt an} der Obergrenze (wie in
unterzeichneten Runden beobachtet) kann sie nicht abbilden.

\hypertarget{trunkierter-erwartungswert}{%
\subsection{15.4 Trunkierter
Erwartungswert}\label{trunkierter-erwartungswert}}

Für überzeichnete Runden ist der bedingte Erwartungswert unterhalb des
Grenzzuschlags erforderlich. Er wird für alle Familien einheitlich über
die Quantilfunktion berechnet -- numerisch stabil auch dort, wo keine
geschlossene Form existiert:

\[ \mathbb{E}\left[b \ \middle|\ b \leq u\right] = \frac{1}{F(u)} \int_{0}^{F(u)} Q(v)\, dv \ \approx\ \frac{1}{n}\sum_{k=1}^{n} Q(v_k), \quad v_k \text{ gleichverteilt in } (0, F(u)) \]

mit \(n = 400\) Stützstellen.

\hypertarget{anpassung-je-runde}{%
\subsection{15.5 Anpassung je Runde}\label{anpassung-je-runde}}

Je Runde \(j\) werden zwei robuste Bedingungen an die Verteilung
gestellt:

\textbf{(a) Mittelwertbedingung.}

\[ \text{unterzeichnet:} \quad \mathbb{E}[b] = \overline{b}_j \]

\[ \text{überzeichnet:} \quad \mathbb{E}\left[b \ \middle|\ b \leq \max_j\right] = \overline{b}_j \]

\textbf{(b) Minimumbedingung.} Das veröffentlichte Minimum der
Zuschlagswerte wird als niedriges Quantil \textbf{aller} Gebote
interpretiert -- das günstigste Gebot gewinnt immer:

\[ Q(\varepsilon_{\min}) = \min_j, \qquad \varepsilon_{\min} = 0{,}02 \]

\begin{quote}
\textbf{Annahme.} \(\varepsilon_{\min} = 2\,\%\) entspricht der
Größenordnung von 40 bis 80 Geboten je Runde. Die Ergebnisse reagieren
auf diese Annahme nur schwach, weil sie den linken Ausläufer betrifft,
während die Gebotsentscheidung am rechten Rand fällt.
\end{quote}

Geschätzt wird über transformierte Parameter, damit die Suche
unbeschränkt laufen kann:

\[ \theta_1 = \mathrm{logit}\left(\mu_{\mathrm{rel}}\right) = \ln\frac{\mu_{\mathrm{rel}}}{1-\mu_{\mathrm{rel}}}, \qquad \theta_2 = \ln \kappa \]

\[ \hat{\theta} = \arg\min_{\theta} \left[ \left(\mathbb{E}_\theta - \overline{b}_j\right)^2 + \left(Q_\theta(\varepsilon_{\min}) - \min_j\right)^2 \right] \]

Verfahren: Trust-Region-Reflective-Kleinstquadrate mit den Schranken
\(\theta_1 \in [-8, 8]\) und \(\kappa \in [1{,}05,\ 800]\), Startwert
\(\mu_{\mathrm{rel}} = 0{,}85\), \(\kappa = 8\).

\textbf{Wettbewerbsquote.} Bei überzeichneten Runden wird sie aus dem
Fit zurückgerechnet. Alle Gebote bis \(p_m\) erhalten einen Zuschlag;
ihr Anteil an allen Geboten ist \(F(p_m)\), und dieser Anteil entspricht
dem Kehrwert der Überzeichnung:

\[ r_j = \frac{1}{F\left(\max_j\right)} \]

Ein Fit-Residuum (Wurzel des mittleren quadratischen Restfehlers beider
Bedingungen) wird je Runde mitgeführt und dient dem Familienvergleich.

\hypertarget{wettbewerbs-link}{%
\subsection{15.6 Wettbewerbs-Link}\label{wettbewerbs-link}}

Über alle Runden hinweg wird der Zusammenhang zwischen
Wettbewerbsintensität und Verteilungsform linear in den transformierten
Größen geschätzt (gewöhnliche Kleinstquadrate):

\[ \mathrm{logit}\left(\mu_{\mathrm{rel},j}\right) = a_L + b_L \ln r_j + \epsilon_j \]

\[ \ln \kappa_j = a_K + b_K \ln r_j + \eta_j \]

Umgekehrt liefert der Link zu einer gegebenen Wettbewerbsquote die
Verteilungsparameter:

\[ \mu_{\mathrm{rel}}(r) = \frac{1}{1 + e^{-\left(a_L + b_L \ln r\right)}}, \qquad \kappa(r) = \max\left(e^{\,a_K + b_K \ln r},\ 1{,}05\right) \]

Die Residuen \(\epsilon_j, \eta_j\) sind das Maß der
Prognoseunsicherheit dieses Zusammenhangs.

\textbf{Lokales Trendmodell.} Für die Rundenprognose werden zusätzlich
die Änderungen der transformierten Parameter zwischen
aufeinanderfolgenden \textbf{Wettbewerbsrunden} ausgewertet:

\[ \Delta_j = \mathrm{logit}\left(\mu_{\mathrm{rel},j}\right) - \mathrm{logit}\left(\mu_{\mathrm{rel},j-1}\right) \]

\[ \mathrm{Drift} = 0, \qquad \mathrm{sd} = \mathrm{clip}\left(\mathrm{StdAbw}\left(\Delta\right),\ 0{,}15,\ 0{,}8\right) \]

\begin{quote}
\textbf{Annahme (Random Walk).} Bei bisher nur drei beobachteten
Rundenänderungen -- davon eine Regimeänderung -- ist die sparsamste
belastbare Annahme ein Drift von null: Die zentrale Prognosewelt
entspricht der letzten Runde, angepasst um Wettbewerbsquote und
Obergrenze. Die \textbf{Streuung} der beobachteten Änderungen geht als
Prognoseunsicherheit ein; sie ist nach oben begrenzt, weil die
Rundenschwankung der Fit-Parameter auch Kalibrierrauschen enthält (aus
Aggregaten ist \(\kappa\) nur grob identifiziert).
\end{quote}

\hypertarget{familienvergleich-und-validierung}{%
\subsection{15.7 Familienvergleich und
Validierung}\label{familienvergleich-und-validierung}}

Die Auswahl der Verteilungsfamilie erfolgt anhand von drei
Validierungskriterien:

\textbf{(1) Massentreue an der Obergrenze.} In unterzeichneten Runden
liegt der Höchstzuschlag an der Obergrenze. Eine Familie muss dort noch
spürbare Masse tragen:

\[ 1 - F\left(0{,}99\,c\right) \geq \frac{1}{50} \]

Ausgewiesen wird der Anteil verletzter Runden.

\textbf{(2) Leave-one-out-Prognose des Grenzzuschlags.} Für jede
überzeichnete Runde wird das Modell \textbf{ohne} diese Runde kalibriert
und der Grenzzuschlag bei gegebener (aus der Runde geschätzter)
Wettbewerbsquote prognostiziert:

\[ \hat{p}_m = Q_{\mu_{\mathrm{rel}}(r),\ \kappa(r)}\left(\min\left(1,\ \frac{1}{r}\right)\right), \qquad \mathrm{RMSE} = \sqrt{\frac{1}{J}\sum_j \left(\hat{p}_{m,j} - p_{m,j}\right)^2} \]

\textbf{(3) Mittleres Fit-Residuum} der beiden Kalibrierbedingungen über
alle Runden.

Zusätzlich wird jede Leave-one-out-Prognose gegen eine \textbf{naive
Basislinie} gestellt -- die Fortschreibung des zuletzt beobachteten
Höchstzuschlags. Diese Referenz bildet den Vergleichsmaßstab für die
Beurteilung, ob das Modell gegenüber einer einfachen Fortschreibung
einen geringeren Prognosefehler erzielt.

\hypertarget{prognose-und-gebotsempfehlung}{%
\subsection{15.8 Prognose und
Gebotsempfehlung}\label{prognose-und-gebotsempfehlung}}

Zwei Modi stehen zur Verfügung.

\hypertarget{modus-letzte-runde-gesetzt}{%
\subsubsection{15.8.1 Modus „letzte Runde
gesetzt``}\label{modus-letzte-runde-gesetzt}}

Die zuletzt beobachtete Ausschreibung gilt als maßgeblich. Verwendet
werden ihre gefittete Verteilung, ihr Grenzzuschlag und ihre
Wettbewerbsquote. Mit dem Zuschlagsanteil

\[ q_c = \min\left(1,\ \frac{1}{r}\right) \]

gilt für ein Gebot \(b\):

\[ \Pr\left[\text{Zuschlag}\right] = \mathrm{clip}\left(\frac{q_c - F(b)}{q_c},\ 0,\ 1\right) \]

Dieser Ausdruck entspricht dem Anteil der Zuschlagswerte der
betrachteten Runde, die oberhalb des eigenen Gebots liegen, und damit
dessen Quantilslage innerhalb der Verteilung. Zur Zielwahrscheinlichkeit
\(z\) folgt das empfohlene Gebot:

\[ b^{*}(z) = Q\left((1-z)\, q_c\right) \]

\hypertarget{modus-prognose-der-nuxe4chsten-runde}{%
\subsubsection{15.8.2 Modus „Prognose der nächsten
Runde``}\label{modus-prognose-der-nuxe4chsten-runde}}

\textbf{Schritt 1: Punktprognosen per Differenzenextrapolation.} Für
Grenzzuschlag und mengengewichteten Mittelwert wird die Zeitreihe der
Wettbewerbsrunden fortgeschrieben. Mit den rekursiven Differenzen

\[ \Delta^{(k)}_t = \Delta^{(k-1)}_t - \Delta^{(k-1)}_{t-1}, \qquad \Delta^{(0)}_t = x_t \]

lautet die Extrapolation der Ordnung \(m\):

\[ D^{(m)} = \Delta^{(m)}_t \]

\[ D^{(k)} = \Delta^{(k)}_t + \lambda_k \cdot D^{(k+1)}, \qquad k = m-1, \dots, 1 \]

\[ \hat{x}_{t+1} = x_t + D^{(1)} \]

Die Dämpfungsparameter \(\lambda_k \in [0,1]\) steuern, wie stark höhere
Ableitungen eingehen: \(\lambda = 1\) übernimmt Trend \textbf{und}
Beschleunigung voll, \(\lambda \rightarrow 0\) nähert sich der linearen
Fortschreibung. Die effektive Ordnung ist durch die Anzahl der
Stützstellen begrenzt; eine Ordnung \(m\) erfordert \(m+1\)
Beobachtungen. Bei nur einer Stützstelle wird der zuletzt beobachtete
Wert fortgeschrieben (Random Walk).

\textbf{Schritt 2: Projektion auf den zulässigen Bereich.}

\[ \hat{p}_m = \mathrm{clip}\left(\hat{x}^{\max}_{t+1},\ 0{,}5,\ c - 0{,}02\right), \qquad \hat{\overline{b}} = \mathrm{clip}\left(\hat{x}^{\,\overline{b}}_{t+1},\ 0{,}3,\ \hat{p}_m - 0{,}05\right) \]

Das Minimum wird aufgrund der fehlenden stabilen Dynamik unverändert
fortgeschrieben (Random Walk) und auf einen Wert unterhalb des
prognostizierten Mittelwerts begrenzt.

\textbf{Schritt 3: Verteilung aus den Punktprognosen.} Methodisch
identisch zum Fit der historischen Runden, mit dem Unterschied, dass die
Mittelwertbedingung stark gewichtet wird (Faktor 8) -- sie ist die vom
Verfahren vorgegebene Größe, während das Minimum als schwächer
gewichtete Nebenbedingung für den linken Ausläufer eingeht:

\[ \hat{\theta} = \arg\min_{\theta}\left[ 64 \left(\mathbb{E}_\theta\left[b \mid b \leq \hat{p}_m\right] - \hat{\overline{b}}\right)^2 + \left(Q_\theta(\varepsilon_{\min}) - \min\right)^2 \right] \]

Die Wettbewerbsquote ist in diesem Modus \textbf{impliziert} und wird
ausgewiesen:

\[ \hat{r} = \frac{1}{F\left(\hat{p}_m\right)} \]

\textbf{Schritt 4: Unsicherheit des Grenzzuschlags.} Um die
Punktprognose wird eine an der Obergrenze und bei 0,5 ct/kWh trunkierte
Normalverteilung gelegt:

\[ p_m \sim \mathcal{N}\left(\hat{p}_m,\ \sigma_{p_m}^2\right) \text{ trunkiert auf } \left[0{,}5,\ c\right] \]

\[ \sigma_{p_m} = \mathrm{clip}\left(\mathrm{StdAbw}\left(\text{Rundenänderungen des Höchstzuschlags}\right),\ 0{,}15,\ 0{,}8\right) \]

Die Trunkierung beschränkt den Grenzzuschlag auf den zulässigen
Wertebereich. Für die numerische Auswertung werden 4.000 Szenarien
simuliert.

\textbf{Schritt 5: Entscheidungsgrößen.} Als Price-Taker gilt: Ein Gebot
\(b\) erhält genau dann einen Zuschlag, wenn der Grenzzuschlag darüber
liegt.

\[ \Pr\left[\text{Zuschlag}(b)\right] = \Pr\left[p_m > b\right] \approx \frac{1}{n}\left|\left\{k : p_m^{(k)} > b\right\}\right| \]

\[ b^{*}(z) = \text{empirisches } (1-z)\text{-Quantil von } \left\{p_m^{(k)}\right\} \]

Aus der Quantildefinition folgt, dass mit zunehmender angestrebter
Zuschlagswahrscheinlichkeit ein niedrigeres Gebot erforderlich ist.

\textbf{Schritt 6: Ziehungen für die Monte-Carlo-Kopplung.} Für die
Simulation werden stochastische Realisationen erfolgreicher
Zuschlagswerte erzeugt:

\[ u \sim \mathcal{U}(0,\ q_c), \qquad b^{\mathrm{basis}} = Q(u) \]

\[ b^{(k)} = \mathrm{clip}\left(b^{\mathrm{basis}} + \left(p_m^{(k)} - \hat{p}_m\right),\ 0,\ c\right) \]

Die zentrale Verteilung wird in jeder Simulationswelt parallel zum
gezogenen Grenzzuschlag verschoben. Die \textbf{Form} bleibt erhalten;
ausschließlich die Lage folgt der Unsicherheit. Im Modus „letzte Runde``
entfällt die Verschiebung.

\hypertarget{deutschland-manuelle-vorgabe}{%
\subsection{15.9 Deutschland: manuelle
Vorgabe}\label{deutschland-manuelle-vorgabe}}

Im deutschen EEG-Marktsystem entfällt das empirische Modell: Der
erwartete Marktprämienzuschlag (anzulegender Wert) wird direkt als Zahl
eingetragen. Die Historie der österreichischen OeMAG-Ausschreibungen ist
für die deutschen Ausschreibungen nicht aussagekräftig. Eine
unmittelbare Übertragung wäre aufgrund der abweichenden
Auktionssystematik methodisch nicht belastbar. Alles Weitere --
Cashflow, Steuern, Kennzahlen -- ist identisch; nur die Herkunft von
\(z\) unterscheidet sich.

\hypertarget{portfolio--und-vergleichsrechnungen}{%
\section{16 Portfolio- und
Vergleichsrechnungen}\label{portfolio--und-vergleichsrechnungen}}

Über die Einzelprojektbewertung hinaus aggregiert die Anwendung über
alle \textbf{aktiven} Projekte. Inaktive Projekte bleiben im
Datenbestand erhalten, werden jedoch bei der Berechnung der
Portfoliokennzahlen nicht berücksichtigt.

\[ P^{\mathrm{ges}} = \sum_{v} P_v, \qquad I^{\mathrm{ges}} = \sum_{v} I_v, \qquad EK^{\mathrm{ges}} = \sum_{v} EK_v \]

\[ \overline{\mathrm{IRR}} = \frac{1}{\left|\mathcal{V}\right|}\sum_{v \in \mathcal{V}} \mathrm{IRR}_v, \qquad \mathcal{V} = \left\{v : \mathrm{IRR}_v \text{ berechenbar}\right\} \]

\begin{quote}
Die Portfolio-Rendite ist ein \textbf{ungewichtetes arithmetisches
Mittel} der Projekt-IRR, keine kapitalgewichtete Portfolio-IRR. Sie
beantwortet die Frage „wie rentabel sind die Projekte im Schnitt``,
nicht „welche Rendite erzielt das eingesetzte Kapital insgesamt``. Für
die zweite Frage müssten die Cashflows aller Projekte zunächst
datumsgenau zusammengeführt und daraus ein gemeinsamer XIRR bestimmt
werden.
\end{quote}

Für die Rendite-Risiko-Landkarte wird je Projekt das spezifische Invest
gebildet:

\[ \text{Invest}_v = \frac{I_v}{P_v} \quad \left[\text{€/kWp}\right] \]

\textbf{Wertbrücke.} Die Brücke über die Gesamtlaufzeit summiert jede
Cashflow-Komponente über alle Jahre und stellt sie als Wasserfall dar:

\[ \sum_t R_t \ -\ \sum_t C_t \ -\ \sum_t Z_t \ -\ \sum_t S_t \ =\ \sum_t \mathrm{CF}^{\mathrm{op}}_t \]

\[ \sum_t \mathrm{CF}^{\mathrm{op}}_t \ -\ I \ +\ D \ -\ \sum_t T_t \ =\ \mathrm{CF}^{\mathrm{kum}}_N \]

Beide Identitäten gelten exakt -- sie sind die Probe darauf, dass die
Cashflow-Aggregation vollständig ist.

\hypertarget{modellannahmen-vereinfachungen-und-geltungsgrenzen}{%
\section{17 Modellannahmen, Vereinfachungen und
Geltungsgrenzen}\label{modellannahmen-vereinfachungen-und-geltungsgrenzen}}

Dieses Kapitel sammelt alle Stellen, an denen das Modell eine
Entscheidung trifft, die über die reine Rechenvorschrift hinausgeht. Wer
Ergebnisse interpretiert, sollte diese Liste kennen.

\hypertarget{zeit-und-perioden}{%
\subsection{17.1 Zeit und Perioden}\label{zeit-und-perioden}}

\begin{longtable}[]{@{}
  >{\raggedright\arraybackslash}p{(\columnwidth - 4\tabcolsep) * \real{0.3333}}
  >{\raggedright\arraybackslash}p{(\columnwidth - 4\tabcolsep) * \real{0.3333}}
  >{\raggedright\arraybackslash}p{(\columnwidth - 4\tabcolsep) * \real{0.3333}}@{}}
\toprule\noalign{}
\begin{minipage}[b]{\linewidth}\raggedright
Nr.
\end{minipage} & \begin{minipage}[b]{\linewidth}\raggedright
Annahme
\end{minipage} & \begin{minipage}[b]{\linewidth}\raggedright
Wirkung
\end{minipage} \\
\midrule\noalign{}
\endhead
\bottomrule\noalign{}
\endlastfoot
A1 & Betriebsperioden sind Kalenderjahre; nur das Anlaufjahr ist
anteilig & Der Sonderfall „Vertragsende am Jahrestag der
Inbetriebnahme`` ist nicht abgebildet \\
A2 & Anteilsfaktor bei 1 gekappt & Schaltjahre erzeugen keinen
366/365-Effekt \\
A3 & Alle Zahlungen eines Jahres wirken zum Jahresende & Unterjährige
Zahlungsströme werden nicht abgebildet; die Diskontierung ist dennoch
taggenau \\
\end{longtable}

\hypertarget{erluxf6se}{%
\subsection{17.2 Erlöse}\label{erluxf6se}}

\begin{longtable}[]{@{}
  >{\raggedright\arraybackslash}p{(\columnwidth - 4\tabcolsep) * \real{0.3333}}
  >{\raggedright\arraybackslash}p{(\columnwidth - 4\tabcolsep) * \real{0.3333}}
  >{\raggedright\arraybackslash}p{(\columnwidth - 4\tabcolsep) * \real{0.3333}}@{}}
\toprule\noalign{}
\begin{minipage}[b]{\linewidth}\raggedright
Nr.
\end{minipage} & \begin{minipage}[b]{\linewidth}\raggedright
Annahme
\end{minipage} & \begin{minipage}[b]{\linewidth}\raggedright
Wirkung
\end{minipage} \\
\midrule\noalign{}
\endhead
\bottomrule\noalign{}
\endlastfoot
A4 & Kurvenwerte außerhalb des Stützbereichs werden geklemmt &
Konservativ gegenüber Extrapolation, unterschätzt aber langfristige
Preistrends \\
A5 & Inflationierung mit dem tatsächlichen Kalenderjahr, auch bei
geklemmtem Realwert & Nominale Fortschreibung des letzten bekannten
Realpreises \\
A6 & Der EAG-Zuschlagswert ist nominal fix & Entspricht der gesetzlichen
Regelung, keine Vereinfachung \\
A7 & Negative Preise wirken über einen jährlichen Mengenanteil, nicht
stundenscharf & Innerjährliche Korrelation von Erzeugung und Preis wird
nicht modelliert \\
A8 & Kein Eigenverbrauch, keine Speicher, keine Regelenergie- oder
PPA-Erlöse & Reines Marktprämien-/Marktverkaufsmodell \\
\end{longtable}

\hypertarget{kosten-und-finanzierung}{%
\subsection{17.3 Kosten und
Finanzierung}\label{kosten-und-finanzierung}}

\begin{longtable}[]{@{}
  >{\raggedright\arraybackslash}p{(\columnwidth - 4\tabcolsep) * \real{0.3333}}
  >{\raggedright\arraybackslash}p{(\columnwidth - 4\tabcolsep) * \real{0.3333}}
  >{\raggedright\arraybackslash}p{(\columnwidth - 4\tabcolsep) * \real{0.3333}}@{}}
\toprule\noalign{}
\begin{minipage}[b]{\linewidth}\raggedright
Nr.
\end{minipage} & \begin{minipage}[b]{\linewidth}\raggedright
Annahme
\end{minipage} & \begin{minipage}[b]{\linewidth}\raggedright
Wirkung
\end{minipage} \\
\midrule\noalign{}
\endhead
\bottomrule\noalign{}
\endlastfoot
A9 & Keine Bauzeitzinsen, keine Zwischenfinanzierung, kein Disagio &
Solche Kosten müssen im CAPEX erfasst werden \\
A10 & Ein Kredit, feste Kondition über die gesamte Laufzeit & Keine
Zinsbindungsfristen, keine Anschlussfinanzierung \\
A11 & Der Anteilsfaktor des Anlaufjahres reduziert nur den Zins, nicht
die Tilgung & Bei Annuität wird im Anlaufjahr modellgemäß ein höherer
Tilgungsanteil angesetzt \\
A12 & Keine Liquiditätsreserve, kein Schuldendienstreservekonto & DSCR
wird ausgewiesen, aber nicht als Nebenbedingung erzwungen \\
A13 & Kein Restwert und keine Rückbaukosten am Ende der Betriebsdauer &
Beides ist gegebenenfalls über CAPEX bzw. eine OPEX-Position
abzubilden \\
A13a & DSCR-Kovenanten wirken nicht auf die Cashflow-Rechnung zurück &
Cash Trap und Equity Cure werden als Prüfung ausgewertet, nicht als
Zahlungsstrom gebucht \\
A13b & Ausgeschüttete Mittel gelten als in voller Höhe rückführbar & Der
Deckungswasserfall beziffert die Obergrenze der Deckung aus eigener
Kraft \\
\end{longtable}

\hypertarget{steuern}{%
\subsection{17.4 Steuern}\label{steuern}}

\begin{longtable}[]{@{}
  >{\raggedright\arraybackslash}p{(\columnwidth - 4\tabcolsep) * \real{0.3333}}
  >{\raggedright\arraybackslash}p{(\columnwidth - 4\tabcolsep) * \real{0.3333}}
  >{\raggedright\arraybackslash}p{(\columnwidth - 4\tabcolsep) * \real{0.3333}}@{}}
\toprule\noalign{}
\begin{minipage}[b]{\linewidth}\raggedright
Nr.
\end{minipage} & \begin{minipage}[b]{\linewidth}\raggedright
Annahme
\end{minipage} & \begin{minipage}[b]{\linewidth}\raggedright
Wirkung
\end{minipage} \\
\midrule\noalign{}
\endhead
\bottomrule\noalign{}
\endlastfoot
A14 & Lineare AfA über die gesamte Investitionssumme & Keine
Komponentenaufteilung, keine degressive Abschreibung, kein
Investitionsfreibetrag \\
A15 & Verlustvortrag zeitlich unbegrenzt, Verrechnungsgrenze je
Gewinnjahr & Entspricht § 8 Abs. 4 Z 2 KStG \\
A16 & Deutsche Gewerbesteuer ohne Verlustvortrag nach § 10a GewStG &
Bewusst, zur Deckungsgleichheit mit dem Referenzmodell \\
A17 & Keine Umsatzsteuer, keine Grundsteuer, keine Kapitalertragsteuer
auf Ausschüttungen & Betrachtung endet auf Ebene der
Projektgesellschaft \\
\end{longtable}

\hypertarget{kennzahlen-und-auswertungen}{%
\subsection{17.5 Kennzahlen und
Auswertungen}\label{kennzahlen-und-auswertungen}}

\begin{longtable}[]{@{}
  >{\raggedright\arraybackslash}p{(\columnwidth - 4\tabcolsep) * \real{0.3333}}
  >{\raggedright\arraybackslash}p{(\columnwidth - 4\tabcolsep) * \real{0.3333}}
  >{\raggedright\arraybackslash}p{(\columnwidth - 4\tabcolsep) * \real{0.3333}}@{}}
\toprule\noalign{}
\begin{minipage}[b]{\linewidth}\raggedright
Nr.
\end{minipage} & \begin{minipage}[b]{\linewidth}\raggedright
Annahme
\end{minipage} & \begin{minipage}[b]{\linewidth}\raggedright
Wirkung
\end{minipage} \\
\midrule\noalign{}
\endhead
\bottomrule\noalign{}
\endlastfoot
A18 & IRR nur bei Vorzeichenwechsel im Cashflow; sonst kein Wert & Kein
Ersatzwert, der als Rendite fehlgelesen werden könnte \\
A19 & Payback undiskontiert, in vollen Jahren & Keine unterjährige
Interpolation \\
A20 & Monte-Carlo-Treiber sind unabhängig normalverteilt & Korrelationen
(z. B. Marktwert-Niveau mit negativen Stunden) sind nicht abgebildet \\
A21 & Fester Startwert des Zufallsgenerators & Reproduzierbar, aber
keine Stichprobenvariation über Läufe hinweg \\
A22 & Portfolio-IRR als ungewichtetes Mittel & Keine kapitalgewichtete
Portfoliorendite \\
\end{longtable}

\hypertarget{auktionsmodell}{%
\subsection{17.6 Auktionsmodell}\label{auktionsmodell}}

\begin{longtable}[]{@{}
  >{\raggedright\arraybackslash}p{(\columnwidth - 4\tabcolsep) * \real{0.3333}}
  >{\raggedright\arraybackslash}p{(\columnwidth - 4\tabcolsep) * \real{0.3333}}
  >{\raggedright\arraybackslash}p{(\columnwidth - 4\tabcolsep) * \real{0.3333}}@{}}
\toprule\noalign{}
\begin{minipage}[b]{\linewidth}\raggedright
Nr.
\end{minipage} & \begin{minipage}[b]{\linewidth}\raggedright
Annahme
\end{minipage} & \begin{minipage}[b]{\linewidth}\raggedright
Wirkung
\end{minipage} \\
\midrule\noalign{}
\endhead
\bottomrule\noalign{}
\endlastfoot
A23 & Price-Taker-Sicht & Das eigene Gebot beeinflusst den Grenzzuschlag
nicht \\
A24 & \(\varepsilon_{\min} = 2\,\%\) für die Minimumbedingung &
Schwacher Einfluss, betrifft nur den linken Ausläufer \\
A25 & Random Walk statt geschätztem Drift & Zentrale Prognose entspricht
der letzten Runde \\
A26 & Wettbewerbsquote überzeichneter Runden aus dem Fit rückgeschätzt &
Das eingereichte Gebotsvolumen wird nicht veröffentlicht \\
A27 & Verteilungsfamilie ist eine Modellwahl & Wird über Massentreue,
Leave-one-out und Fit-Residuum geprüft, nicht gesetzt \\
\end{longtable}

\hypertarget{daten}{%
\subsection{17.7 Daten}\label{daten}}

Die mitgelieferten Preiskurven sind plausible Platzhalter
beziehungsweise Studienauszüge, keine für den Einzelfall validierten
Marktprognosen. Vor einer Investitionsentscheidung sind sie durch
aktuelle, lizenzierte Marktwert-Solar-Kurven zu ersetzen. Das Modell
selbst ist davon unberührt -- es rechnet mit der Kurve, die hinterlegt
ist.

\hypertarget{nachvollziehbarkeit-und-verifikation}{%
\section{18 Nachvollziehbarkeit und
Verifikation}\label{nachvollziehbarkeit-und-verifikation}}

\hypertarget{zuordnung-rechenschritt-zu-code-und-test}{%
\subsection{18.1 Zuordnung Rechenschritt zu Code und
Test}\label{zuordnung-rechenschritt-zu-code-und-test}}

\begin{longtable}[]{@{}
  >{\raggedright\arraybackslash}p{(\columnwidth - 6\tabcolsep) * \real{0.2500}}
  >{\raggedright\arraybackslash}p{(\columnwidth - 6\tabcolsep) * \real{0.2500}}
  >{\raggedright\arraybackslash}p{(\columnwidth - 6\tabcolsep) * \real{0.2500}}
  >{\raggedright\arraybackslash}p{(\columnwidth - 6\tabcolsep) * \real{0.2500}}@{}}
\toprule\noalign{}
\begin{minipage}[b]{\linewidth}\raggedright
Kapitel
\end{minipage} & \begin{minipage}[b]{\linewidth}\raggedright
Rechenschritt
\end{minipage} & \begin{minipage}[b]{\linewidth}\raggedright
Codestelle
\end{minipage} & \begin{minipage}[b]{\linewidth}\raggedright
Test
\end{minipage} \\
\midrule\noalign{}
\endhead
\bottomrule\noalign{}
\endlastfoot
4 & Parameterauflösung &
\nolinkurl{engine/pipeline.py}: \texttt{resolve\_assumptions} &
\nolinkurl{tests/test_pipeline_kpis_io.py} \\
5 & Zeitachse, Anteilsfaktoren & \nolinkurl{engine/timeline.py} &
\nolinkurl{tests/test_timeline_energy_opex.py} \\
6 & Energieertrag & \nolinkurl{engine/energy.py} &
\nolinkurl{tests/test_timeline_energy_opex.py} \\
7 & Erlöse, Marktprämie & \nolinkurl{engine/revenue.py} &
\nolinkurl{tests/test_revenue.py} \\
8 & Betriebskosten & \nolinkurl{engine/opex.py} &
\nolinkurl{tests/test_timeline_energy_opex.py},
\nolinkurl{tests/test_model_optionen.py} \\
9 & Finanzierung & \nolinkurl{engine/financing.py} &
\nolinkurl{tests/test_financing_tax.py} \\
10 & Steuern & \nolinkurl{engine/tax.py} &
\nolinkurl{tests/test_financing_tax.py},
\nolinkurl{tests/test_markt_system.py} \\
11 & Cashflow & \nolinkurl{engine/cashflow.py} &
\nolinkurl{tests/test_pipeline_kpis_io.py} \\
11.5 & DSCR-Kovenanten & \nolinkurl{engine/covenants.py} &
\nolinkurl{tests/test_covenants.py} \\
12 & Kennzahlen, LCOE & \nolinkurl{engine/kpis.py},
\nolinkurl{engine/analytics.py} &
\nolinkurl{tests/test_pipeline_kpis_io.py},
\nolinkurl{tests/test_analytics.py} \\
13 & Beispielrechnung & \nolinkurl{docs/rechenmodell/beispiel.py} &
\nolinkurl{tests/test_dokumentation.py} \\
14 & Sensitivität, Monte Carlo & \nolinkurl{engine/sensitivity.py},
\nolinkurl{engine/analytics.py} & \nolinkurl{tests/test_analytics.py} \\
15 & Auktionsmodell & \nolinkurl{engine/auktion.py} &
\nolinkurl{tests/test_auktion.py} \\
16 & Portfolio, Wertbrücke & \nolinkurl{app/views/overview.py},
\nolinkurl{app/components/charts.py} & \nolinkurl{tests/test_ui_smoke.py} \\
-- & Seitensteuerung, Projektseite & \nolinkurl{app/router.py},
\nolinkurl{app/views/project_page.py} &
\nolinkurl{tests/test_navigation.py} \\
\end{longtable}

\hypertarget{teststrategie}{%
\subsection{18.2 Teststrategie}\label{teststrategie}}

Die Einheitstests rechnen gegen \textbf{handgerechnete Erwartungswerte}
auf einem bewusst vereinfachten Fixture-Projekt: flache Marktwertkurve
von 4 ct/kWh, Inflation aus, keine negativen Stunden, keine
Kosteninflation. Damit ist jeder Erwartungswert im Test selbst
nachrechenbar, und ein Fehler in der Engine kann sich nicht in den
Erwartungswerten verstecken.

Die End-to-End-Tests prüfen zusätzlich strukturelle Eigenschaften, die
unabhängig von konkreten Zahlen gelten müssen: Konsistenz der
Cashflow-Kategorien, Monotonie der NPV-Kurve im Diskontsatz, Monotonie
der Sensitivität im Zuschlagswert, verlustfreie Roundtrips über YAML und
Excel.

Die Fixtures hängen ausdrücklich \textbf{nicht} an den änderbaren
Beispieldaten unter \nolinkurl{data/} -- Nutzer können dort frei editieren,
ohne Tests zu brechen. Die einzige Ausnahme ist
\nolinkurl{tests/test_dokumentation.py}, der explizit prüft, ob die in
Kapitel 13 abgedruckten Zahlen noch zum Beispielprojekt passen.

\hypertarget{symbol-spaltenname-export}{%
\subsection{18.3 Symbol, Spaltenname,
Export}\label{symbol-spaltenname-export}}

Die folgende Tabelle verbindet die Notation dieses Dokuments mit den
Spaltennamen in Cashflow-Tabelle, Excel-Export und Oberfläche.

\begin{longtable}[]{@{}lll@{}}
\toprule\noalign{}
Symbol & Spalte in der Cashflow-Tabelle & Einheit \\
\midrule\noalign{}
\endhead
\bottomrule\noalign{}
\endlastfoot
\(E_t\) & \texttt{produktion\_kwh} & kWh \\
\(m^{\mathrm{real}}_t\) & \texttt{marktwert\_real\_ct\_kwh} & ct/kWh \\
\(m_t\) & \texttt{marktwert\_nominal\_ct\_kwh} & ct/kWh \\
\(s_t\) & \texttt{verguetungssatz\_ct\_kwh} & ct/kWh \\
\(R_t\) & \texttt{erloes\_eur} & € \\
\(R^{\mathrm{markt}}_t\) & \texttt{erloes\_markt\_eur} & € \\
\(R^{\text{Prämie}}_t\) & \texttt{erloes\_praemie\_eur} & € \\
\(C_t\) & \texttt{opex\_gesamt\_eur} & € \\
\(C^{\mathrm{gem}}_t\) & \texttt{gemeindeabgabe\_eur} & € \\
\(C^{\mathrm{dv}}_t\) & \texttt{direktvermarktungskosten\_eur} & € \\
\(C^{(j)}_t\) & eine Spalte je Positionsname (z. B. \texttt{Pacht}) &
€ \\
\(Z_t\) & \texttt{zinsen\_eur} & € \\
\(T_t\) & \texttt{tilgung\_eur} & € \\
\(A_t\) & \texttt{afa\_eur} & € \\
\(G_t\) & \texttt{steuerliches\_ergebnis\_vor\_verlustvortrag\_eur} &
€ \\
\(U_t\) & \texttt{verlustvortrag\_genutzt\_eur} & € \\
\(V_{t+1}\) & \texttt{verlustvortrag\_bestand\_eur} & € \\
\(G^{\mathrm{st}}_t\) & \texttt{steuerliches\_ergebnis\_eur} & € \\
\(S_t\) & \texttt{steuer\_eur} & € \\
\(\mathrm{CF}^{\mathrm{op}}_t\) & \texttt{cf\_operativ\_eur} & € \\
\(\mathrm{CF}^{\mathrm{inv}}_t\) & \texttt{cf\_invest\_eur} & € \\
\(\mathrm{CF}^{\mathrm{fin}}_t\) & \texttt{cf\_finanzierung\_eur} & € \\
\(\mathrm{CF}_t\) & \texttt{cf\_gesamt\_eur} & € \\
\(\mathrm{CF}^{\mathrm{kum}}_t\) & \texttt{cf\_kumuliert\_eur} & € \\
\(\mathrm{DSCR}_t\) & \texttt{dscr} & -- \\
\end{longtable}

\hypertarget{plausibilisierungsfolge-fuxfcr-eine-unabhuxe4ngige-nachrechnung}{%
\subsection{18.4 Plausibilisierungsfolge für eine unabhängige
Nachrechnung}\label{plausibilisierungsfolge-fuxfcr-eine-unabhuxe4ngige-nachrechnung}}

Für eine unabhängige Reproduktion in einer Tabellenkalkulation empfiehlt
sich die folgende Prüfsequenz:

\begin{enumerate}
\def\labelenumi{\arabic{enumi}.}
\tightlist
\item
  \(E_1 = P \cdot h\) -- bei Inbetriebnahme im Januar exakt, ohne
  Faktoren.
\item
  \(m_1\) -- Kurvenwert des Inbetriebnahmejahres mal Inflationsfaktor.
\item
  \(C_1\) -- Summe der spezifischen Positionen mal Leistung plus
  produktionsbasierte Sätze, alle ohne Indexierung.
\item
  \(Z_1 = D \cdot i\) und \(T_1 = \mathrm{Ann} - Z_1\).
\item
  \(G_1 = R_1 - C_1 - Z_1 - A_1\) und daraus \(S_1\).
\item
  \(\mathrm{CF}_1 = R_1 - C_1 - Z_1 - S_1 - T_1\).
\end{enumerate}

Bei Übereinstimmung dieser sechs Größen sind die wesentlichen
Berechnungsbeziehungen des ersten Betriebsjahres verifiziert. Die
Folgeperioden unterscheiden sich durch die in den Kapiteln 5 bis 10
definierten zeitabhängigen Faktoren. Kapitel 13 dokumentiert diese
Prüfung anhand des Beispielparametersatzes.

\hypertarget{reproduzierbarer-dokumentationsaufbau}{%
\section{19 Reproduzierbarer
Dokumentationsaufbau}\label{reproduzierbarer-dokumentationsaufbau}}

\hypertarget{generierung}{%
\subsection{19.1 Generierung}\label{generierung}}

\begin{verbatim}
python docs/rechenmodell/build_pdf.py     # erzeugt Rechenmodell.pdf
python docs/rechenmodell/beispiel.py      # erzeugt die Zahlen zu Kapitel 13
make dokumentation                        # beides in einem Schritt
\end{verbatim}

Die fachliche Quelle ist \nolinkurl{docs/rechenmodell/rechenmodell.md}. Der
Build-Prozess erzeugt daraus zunächst eine LaTeX-Datei und anschließend
das PDF. Erforderlich sind Pandoc, Graphviz, XeLaTeX und
\texttt{latexmk}. Sämtliche mathematischen Ausdrücke werden im PDF als
Vektorsatz ausgegeben.

\hypertarget{pflege-der-fachlichen-quelle}{%
\subsection{19.2 Pflege der fachlichen
Quelle}\label{pflege-der-fachlichen-quelle}}

Inhaltliche Änderungen werden ausschließlich in der Markdown-Quelle
vorgenommen. Die verwendete Formelsyntax muss sowohl von KaTeX für die
Darstellung im Repository als auch von XeLaTeX für den PDF-Satz
unterstützt werden. Mehrzeilige Gleichungen können mit \texttt{aligned},
Fallunterscheidungen mit \texttt{cases} gesetzt werden. Nicht
unterstützte Ausdrücke führen zu einem reproduzierbaren Build-Fehler.

\hypertarget{konsistenz-mit-der-implementierung}{%
\subsection{19.3 Konsistenz mit der
Implementierung}\label{konsistenz-mit-der-implementierung}}

Wird eine Rechenvorschrift in der Engine geändert, sind vier Stellen
nachzuziehen:

\begin{enumerate}
\def\labelenumi{\arabic{enumi}.}
\tightlist
\item
  die betroffene Formel in diesem Dokument,
\item
  gegebenenfalls die Annahmenliste in Kapitel 17,
\item
  die Zahlen in Kapitel 13 (über \texttt{beispiel.py} neu erzeugen),
\item
  das gebaute PDF.
\end{enumerate}

Der Test \nolinkurl{tests/test_dokumentation.py} vergleicht die in Kapitel
13 dokumentierten Zahlen mit den aktuellen Ergebnissen der Engine.
Abweichungen zwischen Dokumentation und Implementierung werden dadurch
im Testlauf sichtbar.

\end{document}
